%----------------------------------------------------------------------------------------
%	PACKAGES AND OTHER DOCUMENT CONFIGURATIONS
%----------------------------------------------------------------------------------------

%----------------------------------------------------------------------------------------
%		Géometrie de la page
%----------------------------------------------------------------------------------------
\documentclass[dvipsnames,french,10pt]{book}

\usepackage[
paperheight=24cm, %hauteur du papier
paperwidth=18cm, %largeur du papier
left=1cm, %marge de gauche
right=1cm, %marge de droite
top=1.5cm, %marge du haut
bottom=1.6cm, %marge du bas
%marginparsep=0pt, %distance entre le texte et les notes de marges 
reversemp, %inverse l'emplacement de la marge
headheight=20.60pt %hauteur du header
%showframe, %permet d'afficher le cadre défini ci-dessus
%bindingoffset=1cm %permet d'ajouter le décalage dû au reliage
]{geometry} %Redéfinition de la taille des pages
\raggedbottom


%----------------------------------------------------------------------------------------
%		Generals
%----------------------------------------------------------------------------------------
%\usepackage{fourier} %!! A changer plus tard !!
\usepackage[scaled]{uarial}
\renewcommand*\familydefault{\sfdefault} %% Only if the base font of the document is to be sans serif
\usepackage{frcursive}
\usepackage[T1]{fontenc} %Accents handling
\usepackage[utf8]{inputenc} % Use UTF-8 encoding
%\usepackage{microtype} % Slightly tweak font spacing for aesthetics
\usepackage[english, francais]{babel} % Language hyphenation and typographical rules
\usepackage{marginnote}

%----------------------------------------------------------------------------------------
%		Graphics
%----------------------------------------------------------------------------------------
\usepackage{xcolor}
\usepackage{graphicx, multicol} % Enhanced support for graphics
\graphicspath{{FIG/}{FIG/NetC_fractions_decimales/}{FIG/NetC_fractions/}{FIG/NetC_les_nombres_entiers/}{FIG/NetC_nombres_decimaux}}
\usepackage{wrapfig}
\usepackage{colortbl}

%\usepackage{xsavebox}
% Il faudrait utiliser xsavebox à l'avenir pour réduire la taille du pdf

%----------------------------------------------------------------------------------------
%		Other packages
%----------------------------------------------------------------------------------------
\usepackage{hyperref}
\hypersetup{
	colorlinks=true, %colorise les liens
	breaklinks=true, %permet le retour à la ligne dans les liens trop longs
	urlcolor= sacado_violet,  %couleur des hyperliens et des QR codes
	linkcolor= sacado_violet, %couleur des liens internes
	plainpages=false  %pour palier à "Bookmark problems can occur when you have duplicate page numbers, for example, if you have a page i and a page 1."
}
\usepackage{tabularx}
\newcolumntype{M}[1]{>{\arraybackslash}m{#1}} %Defines a scalable column type in tabular
\usepackage{booktabs} % Enhances quality of tables
\usepackage{diagbox} % barre en diagonale dans un tableau
\usepackage{multicol}
\usepackage[explicit]{titlesec}
\usepackage{xr}
\usepackage{xspace}
\usepackage{array}
\usepackage{listings}
\usepackage{fancyvrb} %verbatim
\usepackage{stmaryrd}
\usepackage{float}


%----------------------------------------------------------------------------------------
%		Headers and footers
%----------------------------------------------------------------------------------------

\pagestyle{empty}
\usepackage{fancyhdr}
\pagestyle{fancy}
\renewcommand{\headrulewidth}{0pt} % pas de filet sous le header

%----------------------------------------------------------------------------------------
%		Mathematics packages
%----------------------------------------------------------------------------------------
\usepackage{amsthm, amsmath, amssymb, mathrsfs} % Mathematical typesetting
\usepackage{marvosym, wasysym} % More symbols
\usepackage[makeroom]{cancel}
\usepackage{xlop}
\usepackage{pgf,tikz,pgfplots}
\pgfplotsset{compat=1.16}
\usepackage{pgf-pie}
\usetikzlibrary{positioning}
\usetikzlibrary{arrows}
\usepackage{pst-plot,pst-tree,pst-func, pstricks-add,pst-node,pst-text}
%\usepackage{units}
\usepackage{nicefrac}
\usepackage[np]{numprint} %Séparation milliers dans un nombre \np{12345} donne 12 345
\usepackage{multido}
\newcommand{\RNum}[1]{\uppercase\expandafter{\romannumeral #1\relax}}

%----------------------------------------------------------------------------------------
%		New text commands
%----------------------------------------------------------------------------------------
\usepackage{calc}
\usepackage{boites}
 \renewcommand{\arraystretch}{1.6}

%%%%% Pour les imports.
\usepackage{import}

%%%%% Pour faire des boites
\usepackage[tikz]{bclogo}
\usepackage{bclogo}
\usepackage{framed}
\usepackage[skins]{tcolorbox}
\tcbuselibrary{breakable}
\tcbuselibrary{skins}
\usetikzlibrary{quotes,babel,arrows.meta,shadows,decorations.pathmorphing,decorations.markings,patterns}
\usepackage{tikzpagenodes}
\usetikzlibrary{plotmarks}

%%%%% Pour les symboles et les ensembles
%\newcommand{\pp}{\leq}
%\newcommand{\pg}{\geq}
%%\newcommand{\euro}{\eurologo{}}
%\newcommand{\R}{\mathbb{R}}
%\newcommand{\N}{\mathbb{N}}
%\newcommand{\D}{\mathbb{D}}
%\newcommand{\Z}{\mathbb{Z}}
%\newcommand{\Q}{\mathbb{Q}}
%\newcommand{\C}{\mathbb{C}}

%%%%% Pour une double minipage
\newcommand{\mini}[4]{
\begin{minipage}[c]{#1}
#2
\end{minipage}
\hfill
\begin{minipage}[c]{#3}
#4
\end{minipage}
}


%\newcommand\hole[1]{\texttt{\_}}
%\newcommand{\PROP}[1]{\textbf{\underline{#1}}}
%\newcommand{\exercice}{\textcolor{OliveGreen}{Exercice : }}
%\newcommand{\correction}{\textcolor{BurntOrange}{Correction : }}
%\newcommand{\propriete}{\textbf{\underline{Propriété}} : }
%\newcommand{\prop}{\textbf{\underline{Propriété}} : }
%\newcommand{\vocabulaire}{\textbf{\underline{Vocabulaire}} : }
%\newcommand{\voca}{\textbf{\underline{Vocabulaire}} : }

\usepackage{enumitem}
\newlist{todolist}{itemize}{2} %Pour faire des QCM
\setlist[todolist]{label=$\square$} %Pour faire des QCM \begin{todolist} instead of itemize
\renewcommand{\FrenchLabelItem}{\textbullet} %bullet dans les items


%----------------------------------------------------------------------------------------
%		Définition de couleurs pour ...
%----------------------------------------------------------------------------------------

%GEOGEBRA

\definecolor{zzttqq}{rgb}{0.6,0.2,0.} %rouge des polygones
\definecolor{qqqqff}{rgb}{0.,0.,1.}
\definecolor{xdxdff}{rgb}{0.49019607843137253,0.49019607843137253,1.}%bleu
\definecolor{qqwuqq}{rgb}{0.,0.39215686274509803,0.} %vert des angles
\definecolor{ffqqqq}{rgb}{1.,0.,0.} %rouge vif
\definecolor{uuuuuu}{rgb}{0.26666666666666666,0.26666666666666666,0.26666666666666666}
\definecolor{qqzzqq}{rgb}{0.,0.6,0.}
\definecolor{cqcqcq}{rgb}{0.7529411764705882,0.7529411764705882,0.7529411764705882} %gris
\definecolor{qqffqq}{rgb}{0.,1.,0.}
\definecolor{ffdxqq}{rgb}{1.,0.8431372549019608,0.}
\definecolor{ffffff}{rgb}{1.,1.,1.}
\definecolor{ududff}{rgb}{0.30196078431372547,0.30196078431372547,1.}
\definecolor{ffqqff}{rgb}{1.,0.,1.}
\definecolor{ffxfqq}{rgb}{1,0.4980392156862745,0}
\definecolor{ffffqq}{rgb}{1,1,0}
\definecolor{qqttzz}{rgb}{0,0.2,0.6}
\definecolor{qqccqq}{rgb}{0,0.8,0}
\definecolor{qqzzff}{rgb}{0,0.6,1}
\definecolor{qqwwzz}{rgb}{0,0.4,0.6}
\definecolor{eqeqeq}{rgb}{0.8784313725490196,0.8784313725490196,0.8784313725490196}

%SACADO

\definecolor{fond}{HTML}{6818A2}  %couleur des entetes etc.  violet sacado
\definecolor{sacado_purple}{RGB}{94,68,145} %% Violet foncé Sacado
\definecolor{sacado_violet}{RGB}{153,117,224} %% Violet clair Sacado
\definecolor{texte}{HTML}{FFFFFF} % couleur du texte des entetes etc.
\definecolor{sacado_blue}{RGB}{0,129,159} %% Bleu Sacado
\definecolor{sacado_green}{RGB}{59, 157, 38} %% Vert Sacado
\definecolor{sacado_yellow}{RGB}{255,180,0} %% Jaune Sacado
\definecolor{sacado_purple}{RGB}{94,68,145} %% Violet foncé Sacado
\definecolor{sacado_violet}{RGB}{153,117,224} %% Violet clair Sacado
\definecolor{sacado_orange}{HTML}{FF8B69} %% Orange Sacado
\definecolor{sacado_red}{HTML}{A11915} %% Rouge Sacado
\definecolor{sacado_gray}{HTML}{7B7485} %% Gris Sacado
%BOITES 

\definecolor{bleu1}{rgb}{0.54,0.79,0.95} %% Bleu
\definecolor{sapgreen}{rgb}{0.4, 0.49, 0}
\definecolor{dvzfxr}{rgb}{0.7,0.4,0.}
\definecolor{beamer}{rgb}{0.5176470588235295,0.49019607843137253,0.32941176470588235} % couleur beamer
\definecolor{preuveRbeamer}{rgb}{0.8,0.4,0}
\definecolor{sectioncolor}{rgb}{0.24,0.21,0.44}
\definecolor{subsectioncolor}{rgb}{0.1,0.21,0.61}
\definecolor{subsubsectioncolor}{rgb}{0.1,0.21,0.61}
\definecolor{info}{rgb}{0.82,0.62,0}
\definecolor{bleu2}{rgb}{0.38,0.56,0.68}
\definecolor{bleu3}{rgb}{0.24,0.34,0.40}
\definecolor{bleu4}{rgb}{0.12,0.20,0.25}
\definecolor{vert}{rgb}{0.21,0.33,0}
\definecolor{vertS}{rgb}{0.05,0.6,0.42}
\definecolor{red}{rgb}{0.78,0,0}
\definecolor{color5}{rgb}{0,0.4,0.58}
\definecolor{eduscol4B}{rgb}{0.19,0.53,0.64}
\definecolor{eduscol4P}{rgb}{0.62,0.12,0.39}
\definecolor{ill_frame}{HTML}{F0F0F0} %Boite illustration contour
\definecolor{ill_back}{HTML}{F7F7F7}  %Boite illustration background
\definecolor{ill_title}{HTML}{AAAAAA} %Boite illustration titre

%----------------------------------------------------------------------------------------
%		QR codes
%----------------------------------------------------------------------------------------

\usepackage[
%final %Pour la compilation finale
draft %Pour le travail sur les documents
]{qrcode}
\usepackage{fontawesome}
\usepackage{fancyqr}
\FancyQrLoad{flat}
\fancyqrset{
%image=\scalebox{.8}{\includegraphics[scale=1]{sacadoA1.png}},image padding=.5,
l color=sacado_violet,r color=sacado_blue}
\newcommand{\qr}[2]{\centering \fancyqr{https://sacado.xyz/qcm/show_course_from_qrcode/#1}

\vspace{.2cm}

#2} %\qr{id} Pour obtenir un qrcode en indiquant seulement l'id de l'exercice





\newcommand{\miniqr}[3]{
\begin{minipage}[c]{.8\linewidth}
#1
\end{minipage}
\hfill
\fbox{
\begin{minipage}[c]{.18\linewidth}
\begin{center}
\fancyqr{https://sacado.xyz/qcm/show\_course\_from\_qrcode/#2}

\vspace{.2cm}

#3
\end{center}
\end{minipage}
}
}



%practice/frombook/<int:ide>/ pour accéder à un exercice depuis le livre.
%----------------------------------------
%
%   Définitions des environnements "pageCours" et "pageExos"
%
%----------------------------------------

\newcounter{cpt}
\newcounter{exo}
\newcounter{cptr}

\newcommand{\titreChap}{Titre de chapitre à définir}

\renewcommand{\chapter}[3]{
  \stepcounter{chapter}
  \setcounter{exo}{0}
  \setcounter{cpt}{0}
  
%\cleardoublepage  % pour commencer à droite
{\Huge \hfill Chapitre \Roman{chapter}.\\
  \bigskip
  #1\\
  \bigskip {\begin{center}
  \fancyqr[image={\includegraphics[scale=.6]{sacadoA1.png}},image padding=.5,height=5cm]{#2}
  \end{center}}  {\normalsize #3}}
\renewcommand{\titreChap}{#1}

%\ifthenelse{\equal{#2}{}}{}{\par
%  \bigskip\bigskip
%  #2}
\newpage
}

\renewenvironment{leftbar}[1][\hsize]
{%
    \def\FrameCommand
    {%
        {\color{black}\vrule width 0.5pt}%
        \hspace{4pt}%must no space.
        \fboxsep=\FrameSep%\colorbox{yellow}%
    }%
    \MakeFramed{\hsize#1\advance\hsize-\width\FrameRestore}%
}
{\endMakeFramed}


\newcommand{\headerGeneral}[3]{ % intitulé, couleur, qrcode
\begin{tikzpicture}[remember picture,overlay]
\coordinate(NO) at (-2,0);
\coordinate(SW) at (22,1);
\coordinate(titre) at (0,0.2);
%\coordinate(qr) at (16.85,0.);
\shade[left color=#2 , right color=#2 ] (NO) rectangle (SW);
\draw (titre) node[color=texte, anchor=west]{ {\large \bf  #1} \quad\quad \bf {\small \titreChap} };
%\draw (qr) node {\qr{#3}};
\end{tikzpicture}
}

\newenvironment{pageCours}{\lhead{%
\pagecolor{white!100}
\headerGeneral{COURS}{fond!70}{p/1234}
}\begin{leftbar}}{\end{leftbar}\newpage}

\newenvironment{pageAD}{\lhead{%
\pagecolor{sacado_violet!6}
\headerGeneral{APPLICATIONS DIRECTES}{sacado_violet!70}{p/1234}
} }{ \newpage}

\newenvironment{pageParcoursu}{\lhead{%
\pagecolor{sacado_green!6} 
\headerGeneral{PARCOURS 1}{sacado_green}{p/1234}
} }{ \newpage}


\newenvironment{pageParcoursd}{\lhead{%
\pagecolor{sacado_blue!6}
\headerGeneral{PARCOURS 2}{sacado_blue_light}{p/1234}
} }{ \newpage}

\newenvironment{pageParcourst}{\lhead{%
\pagecolor{sacado_red!6}
\headerGeneral{PARCOURS 3}{sacado_red}{p/1234}
} }{ \newpage}

\newenvironment{pageBrouillon}{\lhead{%
\pagecolor{sacado_gray!6}
\headerGeneral{BROUILLON}{sacado_gray}{p/1234}
} }{ \newpage}

\newenvironment{pageRituels}{\lhead{%
\pagecolor{fond!6}
\headerGeneral{RITUELS}{fond!70}{p/1234}
} }{ \newpage}

\newenvironment{pageAuto}{\lhead{%
\pagecolor{sacado_orange!6}
\headerGeneral{AUTOÉVALUATION}{sacado_orange}{p/1234}
} }{ \newpage}

\newenvironment{pageHistoire}{\lhead{%
\pagecolor{olive!6}
\headerGeneral{HISTOIRE}{olive}{p/1234}
} }{ \newpage}



\newenvironment{pageExercices}{\lhead{%
\pagecolor{white!100}
\headerGeneral{ACTIVITÉS}{fond}{p/1234}
}\begin{leftbar}}{\end{leftbar}\newpage}



\fancyfoot[L]{\colorbox{fond!70}{\color{texte}\thepage}}
\fancyfoot[C]{}


\newcommand{\titresec}[2]{\phantom{.}\begin{textblock}{1}[0,1](-1.24,0.25)\colorbox{fond!70}{%
\makebox[0.8cm]{\raisebox{0.05cm}[0.6cm][0.15cm]{\color{texte}\LARGE\bf #1}}}\end{textblock}{\LARGE\bf #2}\\\bigskip}

\renewcommand{\thesection}{\arabic{section}}
\titleformat{\section}{}{%
\hspace{-1.15cm}\colorbox{fond!70}{%
\makebox[0.8cm]{\raisebox{0.05cm}[0.6cm][0.15cm]{\color{texte}\LARGE\bf \thesection}}}}{1em}{\bf \LARGE #1}
  
\renewcommand{\thesubsection}{\arabic{subsection}}
            
\titleformat{\subsection}
{%\begin{textblock}{1}[0,1](-1,0.42) toto
  %\end{textblock}
%\reversemarginpar\marginnote[\rule{0.8cm}{0.8cm}]{}[0pt]  \color{red}\normalfont\Large\bfseries}
}{\hspace{-0.83em}
\colorbox{fond!70}{\makebox[0.6cm]{\raisebox{0cm}[1em][0.2em]\normalfont\large\bfseries\color{texte}\thesubsection}}}{1em}{\bf \large #1}




\makeatletter
\newenvironment{TraitV}[1]{%
% #1 couleur du trait (par défaut CouleurA)
% #2 largeur du trait
% #3 distance entre le trait et le texte
\def\FrameCommand{{\color{#1}\vrule width 2pt}
\hspace{1em}}\MakeFramed {\advance\hsize-\width}}%
{\endMakeFramed}
\makeatother

%----------------------------------------
%
%   Définitions des environnements de Définitions, propriétés...
%
%----------------------------------------

%%%%%%%%%%%%% Définitions
\newenvironment{Def}{%
\medskip \begin{tcolorbox}[widget,colback=sacado_violet!15,colframe=sacado_violet!75!black,
title= \stepcounter{cpt} Définition \thecpt. ]}{%
\end{tcolorbox}\par}


\newenvironment{DefT}[1]{%
\medskip \begin{tcolorbox}[widget,colback=sacado_violet!15,colframe=sacado_violet!75!black,
title= \stepcounter{cpt} Définition \thecpt : #1.]}
{%
\end{tcolorbox}\par}


%%%%%%%%%%%%% Proposition
\newenvironment{Prop}{%
\medskip \begin{tcolorbox}[widget,colback=sacado_blue!15,colframe=sacado_blue!75!black,
title= \stepcounter{cpt} Proposition \thecpt.]}
{%
\end{tcolorbox}\par}


%%%%%%%%%%%%% Propriétés
\newenvironment{Pp}{%
\medskip \begin{tcolorbox}[widget,colback=white!100,colframe=sacado_violet!75!black,
title= \stepcounter{cpt} Propriété \thecpt.]}
{%
\end{tcolorbox}\par}

\newenvironment{PpT}[1]{%
\medskip \begin{tcolorbox}[widget,colback=white!100,colframe=sacado_violet!75!black,
title= \stepcounter{cpt} Propriété \thecpt : #1. ]}
{%
\end{tcolorbox}\par}

\newenvironment{Pps}{%
\medskip \begin{tcolorbox}[widget,colback=white!100,colframe=sacado_violet!75!black,
title= \stepcounter{cpt} Propriétés \thecpt.]}
{%
\end{tcolorbox}\par}


%%%%%%%%%%%%% Conséquence
\newenvironment{Cq}{%
\medskip \begin{tcolorbox}[widget,colback=white,colframe=sacado_blue,
title= \stepcounter{cpt} Conséquence \thecpt.]}
{%
\end{tcolorbox}\par}



%%%%%%%%%%%%% Théorèmes
\newenvironment{ThT}[1]{% théorème avec titre
\medskip \begin{tcolorbox}[widget,colback=white!100,colframe=sacado_violet!75!black,
title= \stepcounter{cpt} Théorème \thecpt : #1.]}
{%
\end{tcolorbox}\par}

\newenvironment{Th}{%
\medskip \begin{tcolorbox}[widget,colback=white!100,colframe=sacado_violet!75!black,
title= \stepcounter{cpt} Théorème \thecpt.]}
{%
\end{tcolorbox}\par}


%%%%%%%%%%%%% Règles
\newenvironment{Reg}{%
\medskip \begin{tcolorbox}[widget,colback=sacado_blue!15,colframe=sacado_blue,
title= \stepcounter{cpt} Règle \thecpt.]}
{%
\end{tcolorbox}\par}

%%%%%%%%%%%%% Représentations
\newenvironment{Rep}{%
\medskip \begin{tcolorbox}[widget,colback=white,colframe=sacado_violet!75!white,
title= \stepcounter{cpt} Représentation \thecpt.]}
{%
\end{tcolorbox}\par}

 
%%%%%%%%%%%%% REMARQUES
\newenvironment{Rq}{%
\medskip \begin{tcolorbox}[widget,colback=sacado_orange!15,colframe=sacado_orange,
title= \stepcounter{cpt} Remarque \thecpt.]}
{%
\end{tcolorbox}\par}

\newenvironment{Rqs}{%
\medskip \begin{tcolorbox}[widget,colback=sacado_orange!15,colframe=sacado_orange,
title= \stepcounter{cpt} Remarques \thecpt.]}
{%
\end{tcolorbox}\par}


%%%%%%%%%%%%% EXEMPLES
\newenvironment{Ex}{%
\medskip \begin{tcolorbox}[widget,colback=white,colframe=sacado_blue_light,
title= \stepcounter{cpt} Exemple \thecpt.]}
{%
\end{tcolorbox}\par}

\newenvironment{Exs}{%
\medskip \begin{tcolorbox}[widget,colback=white!15,colframe=sacado_blue_light,
title= \stepcounter{cpt} Exemples \thecpt.]}
{%
\end{tcolorbox}\par}

 
\newenvironment{ExT}[1]{%
\medskip \begin{tcolorbox}[widget,colback=white,colframe=sacado_blue_light,
title= \stepcounter{cpt} Exemple \thecpt   : #1.]}
{%
\end{tcolorbox}\par}

 
\newenvironment{ExCor}{%
\medskip \begin{tcolorbox}[widget,colback=white,colframe=sacado_blue_light ,
title= \stepcounter{cpt} Exercice corrigé \thecpt.]}
{%
\end{tcolorbox}\par}


%%%%%%%%%%%%% Logique
\newenvironment{Log}{%
\medskip \begin{tcolorbox}[widget,colback=sacado_blue!10,colframe=sacado_blue,
title= \stepcounter{cpt} Logique mathématique \thecpt.]}
{%
\end{tcolorbox}\par}
%%%%%%%%%%%%% Logique avec paramètre
\newenvironment{LogT}[1]{%
\medskip \begin{tcolorbox}[widget,colback=sacado_blue!10,colframe=sacado_blue,
title= \stepcounter{cpt} Logique mathématique \thecpt. #1]}
{%
\end{tcolorbox}\par}

%%%%%%%%%%%%% Preuve
\newenvironment{Pv}[1][]{%
\begin{tcolorbox}[breakable, enhanced,widget, colback=sacado_blue!10!white,boxrule=0pt,frame hidden,
borderline west={1mm}{0mm}{sacado_blue!75}]
\textbf{Preuve#1 : }}
{%
\end{tcolorbox}
\par}


%%%%%%%%%%%%% PreuveROC
\newenvironment{PvR}[1][]{%
\begin{tcolorbox}[breakable, enhanced,widget, colback=sacado_blue!10!white,boxrule=0pt,frame hidden,
borderline west={1mm}{0mm}{sacado_blue!75}]
\textbf{Preuve (ROC)#1 : }}
{%
\end{tcolorbox}
\par}


%%%%%%%%%%%%% DemoExigible
\newenvironment{DemoE}{%
\medskip \begin{tcolorbox}[widget,colback=sacado_blue!10,colframe=sacado_blue,
title= \stepcounter{cpt} Démonstration exigible \thecpt. ]}
{%
\end{tcolorbox}\par}





%%%%%%%%%%%%% Compétences
\newenvironment{Cps}[1][]{%
\vspace{0.4cm}
\begin{tcolorbox}[enhanced, lifted shadow={0mm}{0mm}{0mm}{0mm}%
{black!60!white}, attach boxed title to top left={xshift=5mm, yshift*=-3mm}, coltitle=white, colback=white, boxed title style={colback=sacado_green!100}, colframe=sacado_green!75!black,title=\textbf{Compétences associées#1}]}
{%
\end{tcolorbox}
\par}

%%%%%%%%%%%%% Chapitres connexes
\newenvironment{CCon}[1][]{%
\vspace{0.4cm}
\begin{tcolorbox}[breakable, enhanced,widget, colback=white ,boxrule=0pt,frame hidden,
borderline west={2mm}{0mm}{sacado_violet}]
\textbf{#1}}
{%
\end{tcolorbox}
\par}
%%%%%%%%%%%%% Compétences Collège
\newenvironment{CpsCol}[1][]{%
\vspace{0.4cm}
\begin{tcolorbox}[breakable, enhanced,widget, colback=white ,boxrule=0pt,frame hidden,
borderline west={2mm}{0mm}{sacado_violet}]
\textbf{#1}}
{%
\end{tcolorbox}
\par}


%%%%%%%%%%%%% Illustration
\newenvironment{Ill}{%
\medskip \begin{tcolorbox}[widget,colback=white!15,colframe=sacado_violet!75!black,
title= \stepcounter{cpt} Illustration \thecpt. ]}{%
\end{tcolorbox}\par}

%%%%%%%%%%%%% Rituel
\newenvironment{Rit}{%
\medskip \begin{tcolorbox}[widget,colback=white!15,colframe=sacado_violet!75!black,
title= \stepcounter{cpt} Rituel \thecpt. ]}{%
\end{tcolorbox}\par}


%%%%%%%%%%%%% Méthode
\newenvironment{Mt}{%
\medskip \begin{tcolorbox}[widget,colback=white!15,colframe=sacado_violet!75!black,
title= \stepcounter{cpt} Méthode \thecpt. ]}{%
\end{tcolorbox}\par}

%%%%%%%%%%%%% Méthode
\newenvironment{MtT}[1]{%
\medskip \begin{tcolorbox}[widget,colback=white!15,colframe=sacado_violet!75!black,
title= \stepcounter{cpt} Méthode \thecpt. #1 ]}{%
\end{tcolorbox}\par}


%%%%%%%%%%%%% VocU
\newenvironment{VocU}[1]{%
\medskip \begin{tcolorbox}[widget,colback=white!15,colframe=sacado_violet!75,
title= \stepcounter{cpt} Vocabulaire \thecpt. #1 ]}{%
\end{tcolorbox}\par}


%%%%%%%%%%%%% Notation
\newenvironment{Nt}[1]{%
\medskip \begin{tcolorbox}[widget,colback=white!5,colframe=sacado_red!75,
title= \stepcounter{cptr} Notation \thecptr. #1 ]}{%
\end{tcolorbox}\par}

%%%%%%%%%%%%% Ety
\newenvironment{Ety}[1]{%
\medskip \begin{tcolorbox}[widget,colback=white!15,colframe=sacado_violet!75,
title= \stepcounter{cpt} Étymologie \thecpt. #1 ]}{%
\end{tcolorbox}\par}


%%%%%%%%%%%%% His
\newenvironment{His}[1]{%
\begin{tcolorbox}[right=5mm, enhanced, lifted shadow={0mm}{0mm}{0mm}{0mm}%
{sacado_green_dark!90!white}, attach boxed title to top left={xshift=0.3cm, yshift*=-2mm}, coltitle=sacado_green_dark, colback=sacado_green!10!white, boxed title style={colback=white}, colframe=sacado_green_dark,title= Les mathématiciennes et mathématiciens ]
}{%
\end{tcolorbox}\par}


%%%%%%%%%%%%% Attention
\newenvironment{Att}[1]{%
\medskip \begin{tcolorbox}[widget,colback=sacado_red!5,colframe=sacado_red!95!white,
title= \stepcounter{cpt} Notation \thecpt. #1 ]}{%
\end{tcolorbox}\par}



%%%%%%%%%%%%%%%%%%%%%%%%%%%%%%%%%%%%%%%%%%%%%%%%%%%%%%%%%%%%%%%%%%%%%%%%%%%%%%%%%%%%%%%%%%%%%%%%%%%%%%%%%%%%%%%%%%%%%%
%%%%%%%%%%%%%%%%%%%%%%%%%%%%%%%%%%%%%%%%%%%%%%%%%%%%%%%%%%%%%%%%%%%%%%%%%%%%%%%%%%%%%%%%%%%%%%%%%%%%%%%%%%%%%%%%%%%%%%
%%%%%%%%%%%%%%%%  Exercices                                            %%%%%%%%%%%%%%%%%%%%%%%%%%%%%%%%%%%%%%%%%%%%%%%
%%%%%%%%%%%%%%%%%%%%%%%%%%%%%%%%%%%%%%%%%%%%%%%%%%%%%%%%%%%%%%%%%%%%%%%%%%%%%%%%%%%%%%%%%%%%%%%%%%%%%%%%%%%%%%%%%%%%%%
%%%%%%%%%%%%%%%%%%%%%%%%%%%%%%%%%%%%%%%%%%%%%%%%%%%%%%%%%%%%%%%%%%%%%%%%%%%%%%%%%%%%%%%%%%%%%%%%%%%%%%%%%%%%%%%%%%%%%%
 
 
 
%%%%%%%%%%%%% ExoCad 7 paramètres : Compétences, qrcode , calculatrice, python, scratch, tableur, annales
\newenvironment{ExoCad}[7]{% code avant
\tcbset{top=-0.2cm }
\stepcounter{exo}

\begin{tcolorbox}[right=-5mm, enhanced, lifted shadow={0mm}{0mm}{0mm}{0mm}%
{black!60!white}, attach boxed title to top right={xshift=-0.3cm, yshift*=-2mm}, coltitle=sacado_violet!85!black, colback=white!100!white, boxed title style={colback=white}, colframe=sacado_violet!100!black,title= {\footnotesize  #1}  ]
 
\hspace{-1.3cm} 
\begin{minipage}[t]{0.7cm}

 \begin{tikzpicture}
 	\node[fill=sacado_violet,minimum width=0.7cm]{\textcolor{white}{\bf {\Large \theexo}}};
 \end{tikzpicture}


%%%%%%%%%%%%%%%%%%%%%%%% Condition pour la calculatrice
 \ifthenelse{\equal{#3}{1}}{
 \begin{tikzpicture}
 	\node[minimum width=0.7cm]{\includegraphics[scale=0.5]{MISC/calculator.png} };
 \end{tikzpicture} 
 }{
 \ifthenelse{\equal{#3}{2}}{
 \begin{tikzpicture}
 	\node[minimum width=0.7cm]{\includegraphics[scale=0.5]{MISC/no_calculator.png} };
 \end{tikzpicture} 
 }{}
 }

\end{minipage}
\hfill
\begin{minipage}[t]{17.3cm}
} 
{ 
\end{minipage}%code  après
\hfill
\begin{minipage}[t]{1cm}

\begin{center}
\colorbox{sacado_violet}{\includegraphics[height=1cm]{qrcodes/qrDummy.png}}
\colorbox{white}{ {\footnotesize /b/ABCD} }
\end{center}

\end{minipage}

\end{tcolorbox}
\par 
}

 
%%%%%%%%%%%%% ExoCu 7 paramètres : Compétences, qrcode , calculatrice, python, scratch, tableur, annales
\newenvironment{ExoCu}[7]{% code avant
\tcbset{top=-0.2cm }
\stepcounter{exo}

\begin{tcolorbox}[right=-5mm, enhanced, lifted shadow={0mm}{0mm}{0mm}{0mm}%
{black!60!white}, attach boxed title to top right={xshift=-0.3cm, yshift*=-2mm}, coltitle=sacado_green!85!black, colback=white!100!white, boxed title style={colback=white}, colframe=sacado_green!100!black,title= {\footnotesize  #1}  ]
 
\hspace{-1.3cm} 
\begin{minipage}[t]{0.7cm}

 \begin{tikzpicture}
 	\node[fill=sacado_green,minimum width=0.7cm]{\textcolor{white}{\bf {\Large \theexo}}};
 \end{tikzpicture}


%%%%%%%%%%%%%%%%%%%%%%%% Condition pour la calculatrice
 \ifthenelse{\equal{#3}{1}}{
 \begin{tikzpicture}
 	\node[minimum width=0.7cm]{\includegraphics[scale=0.5]{MISC/calculator.png} };
 \end{tikzpicture} 
 }{
 \ifthenelse{\equal{#3}{2}}{
 \begin{tikzpicture}
 	\node[minimum width=0.7cm]{\includegraphics[scale=0.5]{MISC/no_calculator.png} };
 \end{tikzpicture} 
 }{}
 }

\end{minipage}
\hfill
\begin{minipage}[t]{17.3cm}
} 
{ 
\end{minipage}%code  après
\hfill
\begin{minipage}[t]{1cm}

 
\begin{center}
\colorbox{sacado_green}{\includegraphics[height=1cm]{qrcodes/qrDummy.png}}
\colorbox{white}{ {\footnotesize /b/ABCD}  }
\end{center}
 

\end{minipage}

\end{tcolorbox}
\par
}  


 %%%%%%%%%%%%% ExoCd 7 paramètres : Compétences, qrcode , calculatrice, python, scratch, tableur, annales
\newenvironment{ExoCd}[7]{% code avant
\tcbset{top=-0.2cm }
\stepcounter{exo}

\begin{tcolorbox}[right=-5mm, enhanced, lifted shadow={0mm}{0mm}{0mm}{0mm}%
{black!60!white}, attach boxed title to top right={xshift=-0.3cm, yshift*=-2mm}, coltitle=sacado_blue!85!black, colback=white!100!white, boxed title style={colback=white}, colframe=sacado_blue!100!black,title= {\footnotesize  #1}  ]
 
\hspace{-1.3cm} 
\begin{minipage}[t]{0.7cm}

 \begin{tikzpicture}
 	\node[fill=sacado_blue,minimum width=0.7cm]{\textcolor{white}{\bf {\Large \theexo}}};
 \end{tikzpicture}


%%%%%%%%%%%%%%%%%%%%%%%% Condition pour la calculatrice
 \ifthenelse{\equal{#3}{1}}{
 \begin{tikzpicture}
 	\node[minimum width=0.7cm]{\includegraphics[scale=0.5]{MISC/calculator.png} };
 \end{tikzpicture} 
 }{
 \ifthenelse{\equal{#3}{2}}{
 \begin{tikzpicture}
 	\node[minimum width=0.7cm]{\includegraphics[scale=0.5]{MISC/no_calculator.png} };
 \end{tikzpicture} 
 }{}
 }

\end{minipage}
\hfill
\begin{minipage}[t]{17.3cm}
} 
{ 
\end{minipage}%code  après
\hfill
\begin{minipage}[t]{1cm}

\begin{center}
\colorbox{sacado_blue}{\includegraphics[height=1cm]{qrcodes/qrDummy.png}}
\colorbox{white}{ {\footnotesize /b/ABCD} }
\end{center}

\end{minipage}

\end{tcolorbox}
\par
}  


%%%%%%%%%%%%% ExoCt 7 paramètres : Compétences, qrcode , calculatrice, python, scratch, tableur, annales
\newenvironment{ExoCt}[7]{% code avant
\tcbset{top=-0.2cm }
\stepcounter{exo}

\begin{tcolorbox}[right=-5mm, enhanced, lifted shadow={0mm}{0mm}{0mm}{0mm}%
{black!60!white}, attach boxed title to top right={xshift=-0.3cm, yshift*=-2mm}, coltitle=sacado_red!85!black, colback=white!100!white, boxed title style={colback=white}, colframe=sacado_red!100!black,title= {\footnotesize  #1}  ]
 
\hspace{-1.3cm} 
\begin{minipage}[t]{0.7cm}

 \begin{tikzpicture}
 	\node[fill=sacado_red,minimum width=0.7cm]{\textcolor{white}{\bf {\Large \theexo}}};
 \end{tikzpicture}


%%%%%%%%%%%%%%%%%%%%%%%% Condition pour la calculatrice
 \ifthenelse{\equal{#3}{1}}{
 \begin{tikzpicture}
 	\node[minimum width=0.7cm]{\includegraphics[scale=0.5]{MISC/calculator.png} };
 \end{tikzpicture} 
 }{
 \ifthenelse{\equal{#3}{2}}{
 \begin{tikzpicture}
 	\node[minimum width=0.7cm]{\includegraphics[scale=0.5]{MISC/no_calculator.png} };
 \end{tikzpicture} 
 }{}
 }

\end{minipage}
\hfill
\begin{minipage}[t]{17.3cm}
} 
{ 
\end{minipage}%code  après
\hfill
\begin{minipage}[t]{1cm}

\begin{center}
\colorbox{sacado_red}{\includegraphics[height=1cm]{qrcodes/qrDummy.png}}
\colorbox{white}{ {\footnotesize /b/ABCD} }
\end{center}

\end{minipage}

\end{tcolorbox}
\par
}  

 
 
 

%%%%%%%%%%%%% ExoCu 7 paramètres : compétences , qrcode , calculatrice, python, scratch, tableur, annales
\newenvironment{ExoAuto}[7]{% code avant
\tcbset{top=-0.2cm }
\stepcounter{exo}

\begin{tcolorbox}[right=-5mm, enhanced, lifted shadow={0mm}{0mm}{0mm}{0mm}%
{black!60!white}, attach boxed title to top right={xshift=-0.3cm, yshift*=-2mm}, coltitle=sacado_orange!85!black, colback=white!100!white, boxed title style={colback=white}, colframe=sacado_orange!100!black,title= {\footnotesize  #1}  ]
 
\hspace{-1.3cm} 
\begin{minipage}[t]{0.7cm}

 \begin{tikzpicture}
 	\node[fill=sacado_orange,minimum width=0.7cm]{\textcolor{white}{\bf {\Large \theexo}}};
 \end{tikzpicture}


%%%%%%%%%%%%%%%%%%%%%%%% Condition pour la calculatrice
 \ifthenelse{\equal{#3}{1}}{
 \begin{tikzpicture}
 	\node[minimum width=0.7cm]{\includegraphics[scale=0.5]{MISC/calculator.png} };
 \end{tikzpicture} 
 }{
 \ifthenelse{\equal{#3}{2}}{
 \begin{tikzpicture}
 	\node[minimum width=0.7cm]{\includegraphics[scale=0.5]{MISC/no_calculator.png} };
 \end{tikzpicture} 
 }{}
 }

\end{minipage}
\hfill
\begin{minipage}[t]{17.3cm}
} 
{ 
\end{minipage}%code  après
\hfill
\begin{minipage}[t]{1cm}

\begin{center}
\colorbox{sacado_orange}{\includegraphics[height=1cm]{qrcodes/qrDummy.png}}
\colorbox{white}{  {\footnotesize /b/ABCD}  }
\end{center}

\end{minipage}

\end{tcolorbox}
\par
}  

 

%%%%%%%%%%%%% ExoDec 7 paramètres : compétences , qrcode , calculatrice, python, scratch, tableur, annales
 
\newenvironment{ExoDec}[6]{% code avant
\tcbset{top=-0.2cm }
\stepcounter{exo}

\begin{tcolorbox}[right=-5mm, enhanced, lifted shadow={0mm}{0mm}{0mm}{0mm}%
{black!60!white}, attach boxed title to top right={xshift=-0.3cm, yshift*=-2mm}, coltitle=sacado_violet!85!black, colback=white!100!white, boxed title style={colback=white}, colframe=sacado_violet!100!black,title= {\footnotesize  #1}  ]
 
\hspace{-1.3cm} 
\begin{minipage}[t]{1cm}

 \begin{tikzpicture}
 	\node[fill=sacado_violet,minimum width=0.7cm]{\textcolor{white}{\bf {\Large \theexo}}};
 \end{tikzpicture}


%%%%%%%%%%%%%%%%%%%%%%%% Condition pour la calculatrice
 \ifthenelse{\equal{#3}{1}}{
 \begin{tikzpicture}
 	\node[minimum width=0.7cm]{\includegraphics[scale=0.5]{MISC/calculator.png} };
 \end{tikzpicture} 
 }{
 \ifthenelse{\equal{#3}{2}}{
 \begin{tikzpicture}
 	\node[minimum width=0.7cm]{\includegraphics[scale=0.5]{MISC/no_calculator.png} };
 \end{tikzpicture} 
 }{}
 }

\end{minipage}
\begin{minipage}[t]{17.3cm}
} 
{ 
\end{minipage}
\end{tcolorbox}
\par
}  
%%%%%%%%%%%%%%%%%%%%%%%%%%%%%%%%%%%%%%%%%%%%%%%%%%%%%%%%%%%%%%%%%%%%%%%%%%%%%%%%%%%%%%%%%%%%%%%%%%%%%%%%%%%%%%%%%%%%%%
%%%%%%%%%%%%%%%%%%%%%%%%%%%%%%%%%%%%%%%%%%%%%%%%%%%%%%%%%%%%%%%%%%%%%%%%%%%%%%%%%%%%%%%%%%%%%%%%%%%%%%%%%%%%%%%%%%%%%%
%%%%%%%%%%%%%%%%  Exercices sans qrcode                                %%%%%%%%%%%%%%%%%%%%%%%%%%%%%%%%%%%%%%%%%%%%%%%
%%%%%%%%%%%%%%%%%%%%%%%%%%%%%%%%%%%%%%%%%%%%%%%%%%%%%%%%%%%%%%%%%%%%%%%%%%%%%%%%%%%%%%%%%%%%%%%%%%%%%%%%%%%%%%%%%%%%%%
%%%%%%%%%%%%%%%%%%%%%%%%%%%%%%%%%%%%%%%%%%%%%%%%%%%%%%%%%%%%%%%%%%%%%%%%%%%%%%%%%%%%%%%%%%%%%%%%%%%%%%%%%%%%%%%%%%%%%%
 
 
 
%%%%%%%%%%%%% ExoCad 7 paramètres : Compétences , calculatrice, python, scratch, tableur, annales
\newenvironment{ExoCadN}[6]{% code avant
\tcbset{top=-0.2cm }
\stepcounter{exo}

\begin{tcolorbox}[right=-5mm, enhanced, lifted shadow={0mm}{0mm}{0mm}{0mm}%
{black!60!white}, attach boxed title to top right={xshift=-0.3cm, yshift*=-2mm}, coltitle=sacado_violet!85!black, colback=white!100!white, boxed title style={colback=white}, colframe=sacado_violet!100!black,title= {\footnotesize  #1}  ]
 
\hspace{-1.3cm} 
\begin{minipage}[t]{1cm}

 \begin{tikzpicture}
 	\node[fill=sacado_violet,minimum width=0.7cm]{\textcolor{white}{\bf {\Large \theexo}}};
 \end{tikzpicture}


%%%%%%%%%%%%%%%%%%%%%%%% Condition pour la calculatrice
 \ifthenelse{\equal{#3}{1}}{
 \begin{tikzpicture}
 	\node[minimum width=0.7cm]{\includegraphics[scale=0.5]{MISC/calculator.png} };
 \end{tikzpicture} 
 }{
 \ifthenelse{\equal{#3}{2}}{
 \begin{tikzpicture}
 	\node[minimum width=0.7cm]{\includegraphics[scale=0.5]{MISC/no_calculator.png} };
 \end{tikzpicture} 
 }{}
 }

\end{minipage}
\begin{minipage}[t]{17.3cm}
} 
{ 
\end{minipage}%code  après
\end{tcolorbox}
\par 
}

 
%%%%%%%%%%%%% ExoCu 7 paramètres : Compétences , calculatrice, python, scratch, tableur, annales
\newenvironment{ExoCuN}[6]{% code avant
\tcbset{top=-0.2cm }
\stepcounter{exo}

\begin{tcolorbox}[right=-5mm, enhanced, lifted shadow={0mm}{0mm}{0mm}{0mm}%
{black!60!white}, attach boxed title to top right={xshift=-0.3cm, yshift*=-2mm}, coltitle=sacado_green!85!black, colback=white!100!white, boxed title style={colback=white}, colframe=sacado_green!100!black,title= {\footnotesize  #1}  ]
 
\hspace{-1.3cm} 
\begin{minipage}[t]{1cm}

 \begin{tikzpicture}
 	\node[fill=sacado_green,minimum width=0.7cm]{\textcolor{white}{\bf {\Large \theexo}}};
 \end{tikzpicture}


%%%%%%%%%%%%%%%%%%%%%%%% Condition pour la calculatrice
 \ifthenelse{\equal{#2}{1}}{
 \begin{tikzpicture}
 	\node[minimum width=0.7cm]{\includegraphics[scale=0.5]{MISC/calculator.png} };
 \end{tikzpicture} 
 }{
 \ifthenelse{\equal{#2}{2}}{
 \begin{tikzpicture}
 	\node[minimum width=0.7cm]{\includegraphics[scale=0.5]{MISC/no_calculator.png} };
 \end{tikzpicture} 
 }{}
 }

\end{minipage}
\begin{minipage}[t]{17.3cm}
} 
{ 
\end{minipage}
\end{tcolorbox}
\par
}  


 %%%%%%%%%%%%% ExoCd 6 paramètres : Compétences, calculatrice, python, scratch, tableur, annales
\newenvironment{ExoCdN}[6]{% code avant
\tcbset{top=-0.2cm }
\stepcounter{exo}

\begin{tcolorbox}[right=-5mm, enhanced, lifted shadow={0mm}{0mm}{0mm}{0mm}%
{black!60!white}, attach boxed title to top right={xshift=-0.3cm, yshift*=-2mm}, coltitle=sacado_blue!85!black, colback=white!100!white, boxed title style={colback=white}, colframe=sacado_blue!100!black,title= {\footnotesize  #1}  ]
 
\hspace{-1.3cm} 
\begin{minipage}[t]{1cm}

 \begin{tikzpicture}
 	\node[fill=sacado_blue,minimum width=0.7cm]{\textcolor{white}{\bf {\Large \theexo}}};
 \end{tikzpicture}


%%%%%%%%%%%%%%%%%%%%%%%% Condition pour la calculatrice
 \ifthenelse{\equal{#2}{1}}{
 \begin{tikzpicture}
 	\node[minimum width=0.7cm]{\includegraphics[scale=0.5]{MISC/calculator.png} };
 \end{tikzpicture} 
 }{
 \ifthenelse{\equal{#2}{2}}{
 \begin{tikzpicture}
 	\node[minimum width=0.7cm]{\includegraphics[scale=0.5]{MISC/no_calculator.png} };
 \end{tikzpicture} 
 }{}
 }

\end{minipage}
\begin{minipage}[t]{17.3cm}
} 
{ 
\end{minipage}%code  après
\end{tcolorbox}
\par
}  


%%%%%%%%%%%%% ExoCt 6 paramètres : Compétences,   calculatrice, python, scratch, tableur, annales
\newenvironment{ExoCtN}[6]{% code avant
\tcbset{top=-0.2cm }
\stepcounter{exo}

\begin{tcolorbox}[right=-5mm, enhanced, lifted shadow={0mm}{0mm}{0mm}{0mm}%
{black!60!white}, attach boxed title to top right={xshift=-0.3cm, yshift*=-2mm}, coltitle=sacado_red!85!black, colback=white!100!white, boxed title style={colback=white}, colframe=sacado_red!100!black,title= {\footnotesize  #1}  ]
 
\hspace{-1.3cm} 
\begin{minipage}[t]{1cm}

 \begin{tikzpicture}
 	\node[fill=sacado_red,minimum width=0.7cm]{\textcolor{white}{\bf {\Large \theexo}}};
 \end{tikzpicture}


%%%%%%%%%%%%%%%%%%%%%%%% Condition pour la calculatrice
 \ifthenelse{\equal{#2}{1}}{
 \begin{tikzpicture}
 	\node[minimum width=0.7cm]{\includegraphics[scale=0.5]{MISC/calculator.png} };
 \end{tikzpicture} 
 }{
 \ifthenelse{\equal{#2}{2}}{
 \begin{tikzpicture}
 	\node[minimum width=0.7cm]{\includegraphics[scale=0.5]{MISC/no_calculator.png} };
 \end{tikzpicture} 
 }{}
 }

\end{minipage}
\begin{minipage}[t]{17.3cm}
} 
{ 
\end{minipage}%code  après
\end{tcolorbox}
\par
}  

 
 
 

%%%%%%%%%%%%% ExoCu 7 paramètres : compétences , qrcode , calculatrice, python, scratch, tableur, annales
\newenvironment{ExoAutoN}[6]{% code avant
\tcbset{top=-0.2cm }
\stepcounter{exo}

\begin{tcolorbox}[right=5mm, enhanced, lifted shadow={0mm}{0mm}{0mm}{0mm}%
{black!60!white}, attach boxed title to top right={xshift=-0.3cm, yshift*=-2mm}, coltitle=sacado_orange!85!black, colback=white!100!white, boxed title style={colback=white}, colframe=sacado_orange!100!black,title= {\footnotesize  #1}  ]
 
\hspace{-1.4cm} 
\begin{minipage}[t]{0.7cm}

 \begin{tikzpicture}
 	\node[fill=sacado_orange,minimum width=0.7cm]{\textcolor{white}{\bf {\Large \theexo}}};
 \end{tikzpicture}


%%%%%%%%%%%%%%%%%%%%%%%% Condition pour la calculatrice
 \ifthenelse{\equal{#2}{1}}{
 \begin{tikzpicture}
 	\node[minimum width=0.7cm]{\includegraphics[scale=0.5]{MISC/calculator.png} };
 \end{tikzpicture} 
 }{
 \ifthenelse{\equal{#2}{2}}{
 \begin{tikzpicture}
 	\node[minimum width=0.7cm]{\includegraphics[scale=0.5]{MISC/no_calculator.png} };
 \end{tikzpicture} 
 }{}
 }

\end{minipage}
\hfill
\begin{minipage}[t]{17.3cm}
} 
{ 
\end{minipage}%code  après
\end{tcolorbox}
\par
}  
%%%%%%%%%%%%%%%%%%%%%%%%%%%%%%%%%%%%%%%%%%%%%%%%%%%%%%%%%%%%%%%%%%%%%%%%%%%%%%%%%%%%%%%%%%%%%%%%%%%%%%%%%%%%%%%%%%%%%%
%%%%%%%%%%%%%%%%%%%%%%%%%%%%%%%%%%%%%%%%%%%%%%%%%%%%%%%%%%%%%%%%%%%%%%%%%%%%%%%%%%%%%%%%%%%%%%%%%%%%%%%%%%%%%%%%%%%%%%
%%%%%%%%%%%%%%%%  Exercices   sans contours                            %%%%%%%%%%%%%%%%%%%%%%%%%%%%%%%%%%%%%%%%%%%%%%%
%%%%%%%%%%%%%%%%%%%%%%%%%%%%%%%%%%%%%%%%%%%%%%%%%%%%%%%%%%%%%%%%%%%%%%%%%%%%%%%%%%%%%%%%%%%%%%%%%%%%%%%%%%%%%%%%%%%%%%
%%%%%%%%%%%%%%%%%%%%%%%%%%%%%%%%%%%%%%%%%%%%%%%%%%%%%%%%%%%%%%%%%%%%%%%%%%%%%%%%%%%%%%%%%%%%%%%%%%%%%%%%%%%%%%%%%%%%%%


\newcommand{\Sf}[1]{ \vspace{0.1cm}
{\color{fond}{\Large \textbf{#1}}  } 
} 

\newcommand{\Sfe}[1]{ \vspace{0.1cm}
{\color{sacado_blue}{\Large \textbf{#1}}  } 
} 
% fin de la procédure



%%%%%%%%%%%%% Pointillés ou ligne
\newcommand{\point}[1]{\vspace{0.1cm}\multido{}{#1}{ \dotfill \medskip \endgraf}}
\newcommand{\ligne}[1]{\vspace{0.1cm}\multido{}{#1}{ {\color{cqcqcq}\hrulefill} \medskip \endgraf}}
%----------------------------------------
%
%   Macros et opérateurs
%
%----------------------------------------

\newcommand{\second}{2\up{d}\xspace}
\newcommand{\seconde}{2\up{de}\xspace}
\newcommand{\R}{\mathbb R}
\newcommand{\Rp}{\R_+}
\newcommand{\Rpe}{\R_+^*}
\newcommand{\Rm}{\R_-}
\newcommand{\Rme}{\R_-^*}
\newcommand{\N}{\mathbb N}
\newcommand{\D}{\mathbb D}
\newcommand{\Q}{\mathbb Q}
\newcommand{\Z}{\mathbb Z}
\newcommand{\C}{\mathbb C}
\newcommand{\grs}{\mathfrak S}
\newcommand{\IN}[1]{\llbracket 1,#1\rrbracket}
\newcommand{\card}{\text{Card}\,}
\usepackage{mathrsfs}
\newcommand{\parties}{\mathscr P}
\renewcommand{\epsilon}{\varepsilon}
\newcommand{\rmd}{\text{d}}
\newcommand{\diff}{\mathrm D}
\newcommand{\Id}{{\rm Id}}
\newcommand{\e}{{\rm e}}
\newcommand{\I}{{\rm i}}
\newcommand{\J}{{\rm j}}
\newcommand{\ro}{\circ}
\newcommand{\exu}{\exists\,!\,}
\newcommand{\telq}{\,\, \mid \,\,}
\newcommand{\para}{\raisebox{0.1em}{\text{\footnotesize /\hspace{-0.1em}/}}}   
\newcommand{\vect}[1]{\overrightarrow{#1}}
\newcommand{\scal}[2]{\left(\, #1 \mid #2 \, \right)}
\newcommand{\ortho}[1]{{#1}^\perp}
\newcommand{\veci}{\vec{\text{\it \i}}}
\newcommand{\vecj}{\vec{\text{\it \j}}}
\newcommand{\rep}{$(O;\veci,\vecj,\vec{k})$\xspace}
\newcommand{\Oijk}{$(O, \veci,\vecj,\vec{k})$\xspace}
\newcommand{\rond}{repère orthonormal direct}
\newcommand{\bond}{base orthonormale directe}
\newcommand{\eq}{\Longleftrightarrow}
\newcommand{\implique}{\Longrightarrow}
\newcommand{\noneq}{\ \ \ /\hspace{-1.45em}\eq}
\newcommand{\tend}{\longrightarrow}
\newcommand{\egx}[2]{\underset{#1 \tend #2}=}
\newcommand{\asso}{\longmapsto}
\newcommand{\vers}{\longrightarrow}
\newcommand{\eqn}{~\underset{n \rightarrow \infty}{\sim}~}
\newcommand{\eqx}[2]{~\underset{#1 \rightarrow #2}{\sim}~}

\newcommand{\egn}{~\underset{n \rightarrow \infty}{=}~}
\renewcommand{\descriptionlabel}{\hspace{\labelsep}$\bullet$}
\renewcommand{\bar}{\overline}
\DeclareMathOperator{\ash}{{Argsh}}
\DeclareMathOperator{\cotan}{{cotan}}
\DeclareMathOperator{\ach}{{Argch}}
\DeclareMathOperator{\ath}{{Argth}}
\DeclareMathOperator{\sh}{{sh}}
\DeclareMathOperator{\ch}{{ch}}
\DeclareMathOperator{\Mat}{{Mat}}
\DeclareMathOperator{\Vect}{{Vect}}
\DeclareMathOperator{\trace}{{tr}}
\newcommand{\tr}{{}^{\mathrm t}}
\newcommand{\divi}{~\big|~}
\newcommand{\ndivi}{~\not{\big|}~}
\newcommand{\et}{\wedge}
\newcommand{\ou}{\vee}
\renewcommand{\det}{\operatorname{\text{dét}}}
\DeclareMathOperator{\grad}{{grad}}
\renewcommand{\arcsin}{{\mathop{\mathrm{Arcsin}}}}
\renewcommand{\arccos}{{\mathop{\mathrm{Arccos}}}}
\renewcommand{\arctan}{{\mathop{\mathrm{Arctan}}}}
\renewcommand{\tanh}{{\mathop{\mathrm{th}}}}
\newcommand{\pgcd}{\mathop{\mathrm{pgcd}}}
\newcommand{\ppcm}{\mathop{\mathrm{ppcm}}}
\newcommand{\fonc}[4]{\left\{\begin{tabular}{ccc}$#1$ & $\vers$ & $#2$ \\
$#3$ & $\asso$ & $#4$ \end{tabular}\right.}
\renewcommand{\geq}{\geqslant}
\renewcommand{\leq}{\leqslant}
\renewcommand{\Re}{\text{\rm Re}}
\renewcommand{\Im}{\text{\rm Im}}
\renewcommand{\ker}{\mathop{\mathrm{Ker}}}
\newcommand{\Lin}{\mathcal L}
\newcommand{\GO}{\mathcal O}
\newcommand{\GSO}{\mathcal{SO}}
\newcommand{\GL}{\mathcal{GL}}
\renewcommand{\emptyset}{\varnothing}
%\newcommand{\arc}[1]{\overset{\frown}{#1}}
\newcommand{\rg}{\mathop{\mathrm{rg}}}
\newcommand{\ds}{\displaystyle}
\newcommand{\co}[3]{\begin{pmatrix}#1 \\ #2 \\ #3\end{pmatrix}}
\newcommand{\demi}{\frac 1 2}
\newcommand{\limi}[2]{\underset{#1 \rightarrow #2}\lim}
%\usetikzlibrary{quotes,arrows.meta}
% Pour les pavés droits sans annotations de hLl 
\tikzset{
  nonannotated cuboid/.pic={
    \tikzset{%
      every edge quotes/.append style={midway, auto},
      /cuboid/.cd,
      #1
    }
    \draw [every edge/.append style={pic actions, densely dashed, opacity=.5}, pic actions]
    (0,0,0) coordinate (o) -- ++(-\cubescale*\cubex,0,0) coordinate (a) -- ++(0,-\cubescale*\cubey,0) coordinate (b) edge coordinate [pos=1] (g) ++(0,0,-\cubescale*\cubez)  -- ++(\cubescale*\cubex,0,0) coordinate (c) -- cycle
    (o) -- ++(0,0,-\cubescale*\cubez) coordinate (d) -- ++(0,-\cubescale*\cubey,0) coordinate (e) edge (g) -- (c) -- cycle
    (o) -- (a) -- ++(0,0,-\cubescale*\cubez) coordinate (f) edge (g) -- (d) -- cycle;
%    \path [every edge/.append style={pic actions, |-|}]
%    (b) +(0,-5pt) coordinate (b1) edge ["\cubex \cubeunits"'] (b1 -| c)
%    (b) +(-5pt,0) coordinate (b2) edge ["\cubey \cubeunits"] (b2 |- a)
%    (c) +(3.5pt,-3.5pt) coordinate (c2) edge ["\cubez \cubeunits"'] ([xshift=3.5pt,yshift=-3.5pt]e)
    ;
  },
  /cuboid/.search also={/tikz},
  /cuboid/.cd,
  width/.store in=\cubex,
  height/.store in=\cubey,
  depth/.store in=\cubez,
  units/.store in=\cubeunits,
  scale/.store in=\cubescale,
  width=10,
  height=10,
  depth=10,
  units=cm,
  scale=.1,
}

% Pour les pavés droits avec annotations de hLl
\tikzset{
  annotated cuboid/.pic={
    \tikzset{%
      every edge quotes/.append style={midway, auto},
      /cuboid/.cd,
      #1
    }
    \draw [every edge/.append style={pic actions, densely dashed, opacity=.5}, pic actions]
    (0,0,0) coordinate (o) -- ++(-\cubescale*\cubex,0,0) coordinate (a) -- ++(0,-\cubescale*\cubey,0) coordinate (b) edge coordinate [pos=1] (g) ++(0,0,-\cubescale*\cubez)  -- ++(\cubescale*\cubex,0,0) coordinate (c) -- cycle
    (o) -- ++(0,0,-\cubescale*\cubez) coordinate (d) -- ++(0,-\cubescale*\cubey,0) coordinate (e) edge (g) -- (c) -- cycle
    (o) -- (a) -- ++(0,0,-\cubescale*\cubez) coordinate (f) edge (g) -- (d) -- cycle;
    \path [every edge/.append style={pic actions, |-|}]
    (b) +(0,-5pt) coordinate (b1) edge ["\cubex \cubeunits"'] (b1 -| c)
    (b) +(-5pt,0) coordinate (b2) edge ["\cubey \cubeunits"] (b2 |- a)
    (c) +(3.5pt,-3.5pt) coordinate (c2) edge ["\cubez \cubeunits"'] ([xshift=3.5pt,yshift=-3.5pt]e)
    ;
  },
  /cuboid/.search also={/tikz},
  /cuboid/.cd,
  width/.store in=\cubex,
  height/.store in=\cubey,
  depth/.store in=\cubez,
  units/.store in=\cubeunits,
  scale/.store in=\cubescale,
  width=10,
  height=10,
  depth=10,
  units=cm,
  scale=.1,
}

% Pour les pavés droits avec annotations "hauteur", "Longueur", "largeur".
\tikzset{
  defannotated cuboid/.pic={
    \tikzset{%
      every edge quotes/.append style={midway, auto},
      /cuboid/.cd,
      #1
    }
    \draw [every edge/.append style={pic actions, densely dashed, opacity=.5}, pic actions]
    (0,0,0) coordinate (o) -- ++(-\cubescale*\cubex,0,0) coordinate (a) -- ++(0,-\cubescale*\cubey,0) coordinate (b) edge coordinate [pos=1] (g) ++(0,0,-\cubescale*\cubez)  -- ++(\cubescale*\cubex,0,0) coordinate (c) -- cycle
    (o) -- ++(0,0,-\cubescale*\cubez) coordinate (d) -- ++(0,-\cubescale*\cubey,0) coordinate (e) edge (g) -- (c) -- cycle
    (o) -- (a) -- ++(0,0,-\cubescale*\cubez) coordinate (f) edge (g) -- (d) -- cycle;
    \path [every edge/.append style={pic actions, |-|}]
    (b) +(0,-5pt) coordinate (b1) edge ["Longueur"'] (b1 -| c)
    (b) +(-5pt,0) coordinate (b2) edge ["Hauteur"] (b2 |- a)
    (c) +(3.5pt,-3.5pt) coordinate (c2) edge ["largeur"'] ([xshift=3.5pt,yshift=-3.5pt]e)
    ;
  },
  /cuboid/.search also={/tikz},
  /cuboid/.cd,
  width/.store in=\cubex,
  height/.store in=\cubey,
  depth/.store in=\cubez,
  units/.store in=\cubeunits,
  scale/.store in=\cubescale,
  width=10,
  height=10,
  depth=10,
  units=cm,
  scale=.1,
}

\tikzset{
  defiannotated cuboid/.pic={
    \tikzset{%
      every edge quotes/.append style={midway, auto},
      /cuboid/.cd,
      #1
    }
    \draw [every edge/.append style={pic actions, densely dashed, opacity=.5}, pic actions]
    (0,0,0) coordinate (o) -- ++(-\cubescale*\cubex,0,0) coordinate (a) -- ++(0,-\cubescale*\cubey,0) coordinate (b) edge coordinate [pos=1] (g) ++(0,0,-\cubescale*\cubez)  -- ++(\cubescale*\cubex,0,0) coordinate (c) -- cycle
    (o) -- ++(0,0,-\cubescale*\cubez) coordinate (d) -- ++(0,-\cubescale*\cubey,0) coordinate (e) edge (g) -- (c) -- cycle
    (o) -- (a) -- ++(0,0,-\cubescale*\cubez) coordinate (f) edge (g) -- (d) -- cycle;
    \path [every edge/.append style={pic actions, |-|}]
    (b) +(0,-5pt) coordinate (b1) edge ["c"'] (b1 -| c)
    (b) +(-5pt,0) coordinate (b2) edge ["c"] (b2 |- a)
    (c) +(3.5pt,-3.5pt) coordinate (c2) edge ["c"'] ([xshift=3.5pt,yshift=-3.5pt]e)
    ;
  },
  /cuboid/.search also={/tikz},
  /cuboid/.cd,
  width/.store in=\cubex,
  height/.store in=\cubey,
  depth/.store in=\cubez,
  units/.store in=\cubeunits,
  scale/.store in=\cubescale,
  width=10,
  height=10,
  depth=10,
  units=cm,
  scale=.1,
}

\tikzset{
  defiiannotated cuboid/.pic={
    \tikzset{%
      every edge quotes/.append style={midway, auto},
      /cuboid/.cd,
      #1
    }
    \draw [every edge/.append style={pic actions, densely dashed, opacity=.5}, pic actions]
    (0,0,0) coordinate (o) -- ++(-\cubescale*\cubex,0,0) coordinate (a) -- ++(0,-\cubescale*\cubey,0) coordinate (b) edge coordinate [pos=1] (g) ++(0,0,-\cubescale*\cubez)  -- ++(\cubescale*\cubex,0,0) coordinate (c) -- cycle
    (o) -- ++(0,0,-\cubescale*\cubez) coordinate (d) -- ++(0,-\cubescale*\cubey,0) coordinate (e) edge (g) -- (c) -- cycle
    (o) -- (a) -- ++(0,0,-\cubescale*\cubez) coordinate (f) edge (g) -- (d) -- cycle;
    \path [every edge/.append style={pic actions, |-|}]
    (b) +(0,-5pt) coordinate (b1) edge ["L"'] (b1 -| c)
    (b) +(-5pt,0) coordinate (b2) edge ["h"] (b2 |- a)
    (c) +(3.5pt,-3.5pt) coordinate (c2) edge ["l"'] ([xshift=3.5pt,yshift=-3.5pt]e)
    ;
  },
  /cuboid/.search also={/tikz},
  /cuboid/.cd,
  width/.store in=\cubex,
  height/.store in=\cubey,
  depth/.store in=\cubez,
  units/.store in=\cubeunits,
  scale/.store in=\cubescale,
  width=10,
  height=10,
  depth=10,
  units=cm,
  scale=.1,
}

\tikzset{
  xannotated cuboid/.pic={
    \tikzset{%
      every edge quotes/.append style={midway, auto},
      /cuboid/.cd,
      #1
    }
    \draw [every edge/.append style={pic actions, densely dashed, opacity=.5}, pic actions]
    (0,0,0) coordinate (o) -- ++(-\cubescale*\cubex,0,0) coordinate (a) -- ++(0,-\cubescale*\cubey,0) coordinate (b) edge coordinate [pos=1] (g) ++(0,0,-\cubescale*\cubez)  -- ++(\cubescale*\cubex,0,0) coordinate (c) -- cycle
    (o) -- ++(0,0,-\cubescale*\cubez) coordinate (d) -- ++(0,-\cubescale*\cubey,0) coordinate (e) edge (g) -- (c) -- cycle
    (o) -- (a) -- ++(0,0,-\cubescale*\cubez) coordinate (f) edge (g) -- (d) -- cycle;
    \path [every edge/.append style={pic actions, |-|}]
    (b) +(0,-5pt) coordinate (b1) edge ["1cm"'] (b1 -| c)
    (b) +(-5pt,0) coordinate (b2) edge ["1cm"] (b2 |- a)
    (c) +(3.5pt,-3.5pt) coordinate (c2) edge ["1cm"'] ([xshift=3.5pt,yshift=-3.5pt]e)
    ;
  },
  /cuboid/.search also={/tikz},
  /cuboid/.cd,
  width/.store in=\cubex,
  height/.store in=\cubey,
  depth/.store in=\cubez,
  units/.store in=\cubeunits,
  scale/.store in=\cubescale,
  width=10,
  height=10,
  depth=10,
  units=cm,
  scale=.1,
}

\tikzset{
  yannotated cuboid/.pic={
    \tikzset{%
      every edge quotes/.append style={midway, auto},
      /cuboid/.cd,
      #1
    }
    \draw [every edge/.append style={pic actions, densely dashed, opacity=.5}, pic actions]
    (0,0,0) coordinate (o) -- ++(-\cubescale*\cubex,0,0) coordinate (a) -- ++(0,-\cubescale*\cubey,0) coordinate (b) edge coordinate [pos=1] (g) ++(0,0,-\cubescale*\cubez)  -- ++(\cubescale*\cubex,0,0) coordinate (c) -- cycle
    (o) -- ++(0,0,-\cubescale*\cubez) coordinate (d) -- ++(0,-\cubescale*\cubey,0) coordinate (e) edge (g) -- (c) -- cycle
    (o) -- (a) -- ++(0,0,-\cubescale*\cubez) coordinate (f) edge (g) -- (d) -- cycle;
    \path [every edge/.append style={pic actions, |-|}]
    (b) +(0,-5pt) coordinate (b1) edge ["1mm"'] (b1 -| c)
    (b) +(-5pt,0) coordinate (b2) edge ["1mm"] (b2 |- a)
    (c) +(3.5pt,-3.5pt) coordinate (c2) edge ["1mm"'] ([xshift=3.5pt,yshift=-3.5pt]e)
    ;
  },
  /cuboid/.search also={/tikz},
  /cuboid/.cd,
  width/.store in=\cubex,
  height/.store in=\cubey,
  depth/.store in=\cubez,
  units/.store in=\cubeunits,
  scale/.store in=\cubescale,
  width=10,
  height=10,
  depth=10,
  units=cm,
  scale=.1,
}

%Cubes subdivisés
\newif\ifcuboidshade
\newif\ifcuboidemphedge

\tikzset{
  cuboid/.is family,
  cuboid,
  shiftx/.initial=0,
  shifty/.initial=0,
  dimx/.initial=3,
  dimy/.initial=3,
  dimz/.initial=3,
  scale/.initial=1,
  densityx/.initial=1,
  densityy/.initial=1,
  densityz/.initial=1,
  rotation/.initial=0,
  anglex/.initial=0,
  angley/.initial=90,
  anglez/.initial=225,
  scalex/.initial=1,
  scaley/.initial=1,
  scalez/.initial=0.5,
  front/.style={draw=black,fill=white},
  top/.style={draw=black,fill=white},
  right/.style={draw=black,fill=white},
  shade/.is if=cuboidshade,
  shadecolordark/.initial=black,
  shadecolorlight/.initial=white,
  shadeopacity/.initial=0.15,
  shadesamples/.initial=16,
  emphedge/.is if=cuboidemphedge,
  emphstyle/.style={thick},
}

\newcommand{\tikzcuboidkey}[1]{\pgfkeysvalueof{/tikz/cuboid/#1}}

% Commands
\newcommand{\tikzcuboid}[1]{
    \tikzset{cuboid,#1} % Process Keys passed to command
  \pgfmathsetlengthmacro{\vectorxx}{\tikzcuboidkey{scalex}*cos(\tikzcuboidkey{anglex})*28.452756}
  \pgfmathsetlengthmacro{\vectorxy}{\tikzcuboidkey{scalex}*sin(\tikzcuboidkey{anglex})*28.452756}
  \pgfmathsetlengthmacro{\vectoryx}{\tikzcuboidkey{scaley}*cos(\tikzcuboidkey{angley})*28.452756}
  \pgfmathsetlengthmacro{\vectoryy}{\tikzcuboidkey{scaley}*sin(\tikzcuboidkey{angley})*28.452756}
  \pgfmathsetlengthmacro{\vectorzx}{\tikzcuboidkey{scalez}*cos(\tikzcuboidkey{anglez})*28.452756}
  \pgfmathsetlengthmacro{\vectorzy}{\tikzcuboidkey{scalez}*sin(\tikzcuboidkey{anglez})*28.452756}
  \begin{scope}[xshift=\tikzcuboidkey{shiftx}, yshift=\tikzcuboidkey{shifty}, scale=\tikzcuboidkey{scale}, rotate=\tikzcuboidkey{rotation}, x={(\vectorxx,\vectorxy)}, y={(\vectoryx,\vectoryy)}, z={(\vectorzx,\vectorzy)}]
    \pgfmathsetmacro{\steppingx}{1/\tikzcuboidkey{densityx}}
  \pgfmathsetmacro{\steppingy}{1/\tikzcuboidkey{densityy}}
  \pgfmathsetmacro{\steppingz}{1/\tikzcuboidkey{densityz}}
  \newcommand{\dimx}{\tikzcuboidkey{dimx}}
  \newcommand{\dimy}{\tikzcuboidkey{dimy}}
  \newcommand{\dimz}{\tikzcuboidkey{dimz}}
  \pgfmathsetmacro{\secondx}{2*\steppingx}
  \pgfmathsetmacro{\secondy}{2*\steppingy}
  \pgfmathsetmacro{\secondz}{2*\steppingz}
  \foreach \x in {\steppingx,\secondx,...,\dimx}
  { \foreach \y in {\steppingy,\secondy,...,\dimy}
    {   \pgfmathsetmacro{\lowx}{(\x-\steppingx)}
      \pgfmathsetmacro{\lowy}{(\y-\steppingy)}
      \filldraw[cuboid/front] (\lowx,\lowy,\dimz) -- (\lowx,\y,\dimz) -- (\x,\y,\dimz) -- (\x,\lowy,\dimz) -- cycle;
    }
    }
  \foreach \x in {\steppingx,\secondx,...,\dimx}
  { \foreach \z in {\steppingz,\secondz,...,\dimz}
    {   \pgfmathsetmacro{\lowx}{(\x-\steppingx)}
      \pgfmathsetmacro{\lowz}{(\z-\steppingz)}
      \filldraw[cuboid/top] (\lowx,\dimy,\lowz) -- (\lowx,\dimy,\z) -- (\x,\dimy,\z) -- (\x,\dimy,\lowz) -- cycle;
        }
    }
    \foreach \y in {\steppingy,\secondy,...,\dimy}
  { \foreach \z in {\steppingz,\secondz,...,\dimz}
    {   \pgfmathsetmacro{\lowy}{(\y-\steppingy)}
      \pgfmathsetmacro{\lowz}{(\z-\steppingz)}
      \filldraw[cuboid/right] (\dimx,\lowy,\lowz) -- (\dimx,\lowy,\z) -- (\dimx,\y,\z) -- (\dimx,\y,\lowz) -- cycle;
    }
  }
  \ifcuboidemphedge
    \draw[cuboid/emphstyle] (0,\dimy,0) -- (\dimx,\dimy,0) -- (\dimx,\dimy,\dimz) -- (0,\dimy,\dimz) -- cycle;%
    \draw[cuboid/emphstyle] (0,\dimy,\dimz) -- (0,0,\dimz) -- (\dimx,0,\dimz) -- (\dimx,\dimy,\dimz);%
    \draw[cuboid/emphstyle] (\dimx,\dimy,0) -- (\dimx,0,0) -- (\dimx,0,\dimz);%
    \fi

    \ifcuboidshade
    \pgfmathsetmacro{\cstepx}{\dimx/\tikzcuboidkey{shadesamples}}
    \pgfmathsetmacro{\cstepy}{\dimy/\tikzcuboidkey{shadesamples}}
    \pgfmathsetmacro{\cstepz}{\dimz/\tikzcuboidkey{shadesamples}}
    \foreach \s in {1,...,\tikzcuboidkey{shadesamples}}
    {   \pgfmathsetmacro{\lows}{\s-1}
        \pgfmathsetmacro{\cpercent}{(\lows)/(\tikzcuboidkey{shadesamples}-1)*100}
        \fill[opacity=\tikzcuboidkey{shadeopacity},color=\tikzcuboidkey{shadecolorlight}!\cpercent!\tikzcuboidkey{shadecolordark}] (0,\s*\cstepy,\dimz) -- (\s*\cstepx,\s*\cstepy,\dimz) -- (\s*\cstepx,0,\dimz) -- (\lows*\cstepx,0,\dimz) -- (\lows*\cstepx,\lows*\cstepy,\dimz) -- (0,\lows*\cstepy,\dimz) -- cycle;
        \fill[opacity=\tikzcuboidkey{shadeopacity},color=\tikzcuboidkey{shadecolorlight}!\cpercent!\tikzcuboidkey{shadecolordark}] (0,\dimy,\s*\cstepz) -- (\s*\cstepx,\dimy,\s*\cstepz) -- (\s*\cstepx,\dimy,0) -- (\lows*\cstepx,\dimy,0) -- (\lows*\cstepx,\dimy,\lows*\cstepz) -- (0,\dimy,\lows*\cstepz) -- cycle;
        \fill[opacity=\tikzcuboidkey{shadeopacity},color=\tikzcuboidkey{shadecolorlight}!\cpercent!\tikzcuboidkey{shadecolordark}] (\dimx,0,\s*\cstepz) -- (\dimx,\s*\cstepy,\s*\cstepz) -- (\dimx,\s*\cstepy,0) -- (\dimx,\lows*\cstepy,0) -- (\dimx,\lows*\cstepy,\lows*\cstepz) -- (\dimx,0,\lows*\cstepz) -- cycle;
    }
    \fi 

  \end{scope}
}

%\makeatother




\begin{document}

\section{Comparer, estimer mesurer des grandeurs géométriques avec des nombres entier, décimaux : longueur, aire, volume, angle, durée.}

\subsection{Connaître la formule de la longueur d'un cercle et l'utiliser.}

\ExoCu{Calculer.}

On considère le cercle de centre O et de rayon 50 cm. Détermine mentalement une valeur approchée de son périmètre à l'unité près.
 
\begin{minipage}{0.25\linewidth}

\definecolor{uuuuuu}{rgb}{0.26666666666666666,0.26666666666666666,0.26666666666666666}
\definecolor{qqqqff}{rgb}{0.,0.,1.}
\definecolor{ffqqqq}{rgb}{1.,0.,0.}
\definecolor{ududff}{rgb}{0.30196078431372547,0.30196078431372547,1.}
\begin{tikzpicture}[line cap=round,line join=round,>=triangle 45,x=0.5928853754940714cm,y=0.5928853754940714cm]
\clip(-0.4,1.64) rectangle (4.66,6.38);
\draw [line width=1.pt,color=ffqqqq] (2.,4.) circle (1.3257320024702317cm);
\draw [line width=1.pt,color=qqqqff] (2.,4.)-- (4.,3.);
\begin{scriptsize}
\draw [color=ududff] (2.,4.)-- ++(-1.5pt,0 pt) -- ++(3.0pt,0 pt) ++(-1.5pt,-1.5pt) -- ++(0 pt,3.0pt);
\draw[color=ududff] (1.66,4.49) node {$O$};
\end{scriptsize}
\draw [color=ududff](2.5,4.3) node[anchor=north west] {$50 ~cm$};
\end{tikzpicture}
\end{minipage}
\begin{minipage}{0.7\linewidth}

\ligne{6}

\end{minipage}

\ExoCu{Calculer.}

On considère le cercle de centre O et de rayon 30 cm. Détermine son périmètre.

\begin{minipage}{0.25\linewidth}
 
\definecolor{uuuuuu}{rgb}{0.26666666666666666,0.26666666666666666,0.26666666666666666}
\definecolor{qqqqff}{rgb}{0.,0.,1.}
\definecolor{ffqqqq}{rgb}{1.,0.,0.}
\definecolor{ududff}{rgb}{0.30196078431372547,0.30196078431372547,1.}
\begin{tikzpicture}[line cap=round,line join=round,>=triangle 45,x=0.5928853754940714cm,y=0.5928853754940714cm]
\clip(-0.4,1.64) rectangle (4.66,6.38);
\draw [line width=1.pt,color=ffqqqq] (2.,4.) circle (1.3257320024702317cm);
\draw [line width=1.pt,color=qqqqff] (2.,4.)-- (4.,3.);
\begin{scriptsize}
\draw [color=ududff] (2.,4.)-- ++(-1.5pt,0 pt) -- ++(3.0pt,0 pt) ++(-1.5pt,-1.5pt) -- ++(0 pt,3.0pt);
\draw[color=ududff] (1.66,4.49) node {$O$};
\end{scriptsize}
\draw [color=ududff](2.5,4.3) node[anchor=north west] {$50 ~cm$};
\end{tikzpicture}

\end{minipage}
\begin{minipage}{0.7\linewidth}

\ligne{6}

\end{minipage}





\ExoCd{Calculer.}

On considère le cercle de centre $\Omega$ et de rayon $20$ cm. Détermine mentalement une valeur approchée de son périmètre au dixième près.

\begin{minipage}{0.25\linewidth}
 
\definecolor{uuuuuu}{rgb}{0.26666666666666666,0.26666666666666666,0.26666666666666666}
\definecolor{qqqqff}{rgb}{0.,0.,1.}
\definecolor{ffqqqq}{rgb}{1.,0.,0.}
\definecolor{ududff}{rgb}{0.30196078431372547,0.30196078431372547,1.}
\begin{tikzpicture}[line cap=round,line join=round,>=triangle 45,x=0.5928853754940714cm,y=0.5928853754940714cm]
\clip(-0.4,1.64) rectangle (4.66,6.38);
\draw [line width=1.pt,color=ffqqqq] (2.,4.) circle (1.3257320024702317cm);
\draw [line width=1.pt,color=qqqqff] (2.,4.)-- (4.,3.);
\draw [color=ududff](2.6,4.3) node[anchor=north west] {$20 ~cm$};
\begin{scriptsize}
\draw [color=ududff] (2.,4.)-- ++(-1.5pt,0 pt) -- ++(3.0pt,0 pt) ++(-1.5pt,-1.5pt) -- ++(0 pt,3.0pt);
\draw[color=ududff] (1.66,4.39) node {$\Omega$};
\draw [color=ududff] (4.,3.)-- ++(-1.5pt,0 pt) -- ++(3.0pt,0 pt) ++(-1.5pt,-1.5pt) -- ++(0 pt,3.0pt);
\end{scriptsize}
\end{tikzpicture}


\end{minipage}
\begin{minipage}{0.7\linewidth}

\ligne{6}

\end{minipage}


\ExoCu{Calculer.}

On considère le cercle de centre $\Phi$ et de rayon $24,6$ cm. Détermine son périmètre.

\begin{minipage}{0.25\linewidth}

\definecolor{uuuuuu}{rgb}{0.26666666666666666,0.26666666666666666,0.26666666666666666}
\definecolor{qqqqff}{rgb}{0.,0.,1.}
\definecolor{ffqqqq}{rgb}{1.,0.,0.}
\definecolor{ududff}{rgb}{0.30196078431372547,0.30196078431372547,1.}
\begin{tikzpicture}[line cap=round,line join=round,>=triangle 45,x=0.5928853754940714cm,y=0.5928853754940714cm]
\clip(-0.4,1.64) rectangle (4.66,6.38);
\draw [line width=1.pt,color=ffqqqq] (2.,4.) circle (1.3257320024702317cm);
\draw [line width=1.pt,color=qqqqff] (2.,4.)-- (4.,3.);
\begin{scriptsize}
\draw [color=ududff] (2.,4.)-- ++(-1.5pt,0 pt) -- ++(3.0pt,0 pt) ++(-1.5pt,-1.5pt) -- ++(0 pt,3.0pt);
\draw[color=ududff] (1.76,4.39) node {$\Phi$};
\draw [color=ududff] (4.,3.)-- ++(-1.5pt,0 pt) -- ++(3.0pt,0 pt) ++(-1.5pt,-1.5pt) -- ++(0 pt,3.0pt);
\end{scriptsize}
\end{tikzpicture}

\end{minipage}
\begin{minipage}{0.7\linewidth}

\ligne{6}

\end{minipage}



\ExoCu{Calculer. Représenter.}

On considère le cercle de centre $A$ et de diamètre $6$ cm. Détermine son périmètre.
 

\ExoCd{Calculer. Représenter.}

On considère le cercle de centre $A$ et de diamètre $15$ cm. Détermine son périmètre.
 

\ExoCd{Calculer.}

Détermine le périmètre de la figure violette composée d'arcs de cercle.

\begin{minipage}{0.25\linewidth}

\definecolor{xfqqff}{rgb}{0.4980392156862745,0.,1.}
\definecolor{eqeqeq}{rgb}{0.8784313725490196,0.8784313725490196,0.8784313725490196}
\definecolor{uuuuuu}{rgb}{0.26666666666666666,0.26666666666666666,0.26666666666666666}
\definecolor{wwqqcc}{rgb}{0.4,0.,0.8}
\definecolor{aqaqaq}{rgb}{0.6274509803921569,0.6274509803921569,0.6274509803921569}
\definecolor{cqcqcq}{rgb}{0.7529411764705882,0.7529411764705882,0.7529411764705882}
\definecolor{ududff}{rgb}{0.30196078431372547,0.30196078431372547,1.}
\begin{tikzpicture}[line cap=round,line join=round,>=triangle 45,x=0.7220216606498198cm,y=0.7220216606498198cm]
\clip(-1.68,1.6) rectangle (3.86,6.86);
\fill[line width=1.pt,color=eqeqeq,fill=eqeqeq,fill opacity=0.10000000149011612] (-1.4,1.78) -- (3.48,1.78) -- (3.48,6.66) -- (-1.4,6.66) -- cycle;
\draw [line width=1.4pt,color=cqcqcq] (1.,4.) circle (1.489676098115096cm);
\draw [line width=1.pt,color=aqaqaq] (1.,4.)-- (2.98,3.42);
\draw [line width=1.pt,color=aqaqaq] (1.9757503004512438,3.8392166632550544) -- (1.9082823352408018,3.6088949889159623);
\draw [line width=1.pt,color=aqaqaq] (2.0717176647591993,3.8111050110840377) -- (2.0042496995487573,3.5807833367449455);
\draw (1.92,4.52) node[anchor=north west] {$8$};
\draw [shift={(1.99,3.71)},line width=1.pt,color=wwqqcc]  plot[domain=2.856635303339639:5.998227956929432,variable=\t]({1.*1.0316006979447037*cos(\t r)+0.*1.0316006979447037*sin(\t r)},{0.*1.0316006979447037*cos(\t r)+1.*1.0316006979447037*sin(\t r)});
\draw [shift={(0.01,4.29)},line width=1.pt,color=wwqqcc]  plot[domain=-0.28495735025015456:2.856635303339639,variable=\t]({1.*1.0316006979447037*cos(\t r)+0.*1.0316006979447037*sin(\t r)},{0.*1.0316006979447037*cos(\t r)+1.*1.0316006979447037*sin(\t r)});
\draw [line width=1.pt,color=aqaqaq] (-0.98,4.58)-- (1.,4.);
\draw [line width=1.pt,color=aqaqaq] (-0.004249699548757347,4.419216663255054) -- (-0.07171766475919808,4.188894988915963);
\draw [line width=1.pt,color=aqaqaq] (0.09171766475919807,4.391105011084037) -- (0.02424969954875734,4.160783336744945);
\draw [line width=1.pt,color=eqeqeq] (-1.4,1.78)-- (3.48,1.78);
\draw [line width=1.pt,color=eqeqeq] (3.48,1.78)-- (3.48,6.66);
\draw [line width=1.pt,color=eqeqeq] (3.48,6.66)-- (-1.4,6.66);
\draw [line width=1.pt,color=eqeqeq] (-1.4,6.66)-- (-1.4,1.78);
\draw [shift={(1.,4.)},line width=1.pt,color=xfqqff]  plot[domain=-0.28495735025015456:2.856635303339639,variable=\t]({1.*2.0632013958894073*cos(\t r)+0.*2.0632013958894073*sin(\t r)},{0.*2.0632013958894073*cos(\t r)+1.*2.0632013958894073*sin(\t r)});
\begin{scriptsize}
\draw [fill=ududff] (1.,4.) circle (2.5pt);
\draw[color=ududff] (1.14,4.37) node {$A$};
\draw [fill=ududff] (2.98,3.42) circle (2.0pt);
\draw[color=ududff] (3.12,3.75) node {$B$};
\draw[color=wwqqcc] (1.88,3.03) node {$d$};
\draw [fill=uuuuuu] (-0.98,4.58) circle (2.0pt);
\draw[color=uuuuuu] (-1.26,4.85) node {$C$};
\draw[color=wwqqcc] (0.48,5.59) node {$e$};
\draw[color=aqaqaq] (-0.02,4.17) node {$h$};
\end{scriptsize}
\end{tikzpicture}

\end{minipage}
\begin{minipage}{0.7\linewidth}

\ligne{6}

\end{minipage}




\subsection{ Réaliser des conversions nécessitant deux étapes de traitement. (Transformer des heures en semaines, jours et heures ; transformer des secondes en heures, minutes, secondes).}

\ExoCu{Calculer.}

Eva et sa famille sont parties à 20h45 de leur maison. Elles sont arrivées à 22h55 sur leur lieu de vacances. Quelle est la durée de leur voyage ?

\point{3}

\ExoCu{Calculer.}

Convertis 
\begin{enumerate}
\item $120$ minutes = $\ldots\ldots$ minutes 
\item $3$ heures en minutes = $\ldots\ldots$ minutes 
\item $240$ secondes en minutes = $\ldots\ldots$ minutes 
\end{enumerate}

\ExoCu{Calculer.}

Convertis

\begin{enumerate}
\item $21$ h  $12$ min = $\ldots\ldots$ min
\item $12$ min  $55$ s = $\ldots\ldots$ s
\item $17$ min  $9$ min = $\ldots\ldots$ min
\end{enumerate}



\ExoCu{Calculer.}

Convertis

\begin{enumerate}
\item $92$ minutes = $\ldots$ h  $\ldots$ min.
\item $123$ secondes = $\ldots$ min  $\ldots$ s. 
\item $432$ minutes = $\ldots$ h  $\ldots$ min.
\end{enumerate}



\ExoCd{Calculer.}

Convertis

\begin{enumerate}
\item $888$ minutes = $\ldots$ h  $\ldots$ min.
\item $3823$ secondes = $\ldots$ h $\ldots$ min  $\ldots$ s. 
\item $4532$ minutes = $\ldots$ h  $\ldots$ min  $\ldots$ s. 
\end{enumerate}



\ExoCd{Calculer.}

Convertis

\begin{enumerate}
\item $90$ minutes en heures = $\ldots\ldots$ h.
\item $366$ secondes en minutes = $\ldots\ldots$ h.
\end{enumerate}







\subsection{ Utiliser les multiples et sous-multiples du m² et les relations qui les lient.}







\subsection{ Calculer l’aire d’un triangle à l’aide de la formule }

\ExoCu{Calculer.}




\begin{minipage}{0.7\linewidth}
Le triangle ABC est rectangle en A et AB = $10$ cm et AC = $6$ cm. Calcule l'aire du triangle ABC. 

Ta réponse : $ {A}_{ABC}=\ldots\ldots ~ \text{cm}^2$
\end{minipage}
\begin{minipage}{0.5\linewidth}
\definecolor{eqeqeq}{rgb}{0.8784313725490196,0.8784313725490196,0.8784313725490196}
\definecolor{qqwuqq}{rgb}{0.,0.39215686274509803,0.}
\definecolor{qqqqff}{rgb}{0.,0.,1.}
\definecolor{ududff}{rgb}{0.30196078431372547,0.30196078431372547,1.}
\begin{tikzpicture}[line cap=round,line join=round,>=triangle 45,x=0.6309148264984229cm,y=0.6756756756756758cm]
\clip(-2.2,1.8) rectangle (4.14,6.24);
\draw[line width=1.pt,color=qqwuqq,fill=qqwuqq,fill opacity=0.1] (-0.5757359312880717,3.) -- (-0.5757359312880715,3.4242640687119286) -- (-1.,3.4242640687119286) -- (-1.,3.) -- cycle; 
\fill[line width=1.pt,color=eqeqeq,fill=eqeqeq,fill opacity=0.1] (-2.,2.) -- (4.,2.) -- (4.,6.) -- (-2.,6.) -- cycle;
\draw (0.62,3.16) node[anchor=north west] {$8$};
\draw [line width=1.pt,color=qqqqff] (-1.,5.)-- (-1.,3.);
\draw [line width=1.pt,color=qqqqff] (-1.,3.)-- (3.,3.);
\draw [line width=1.pt,color=qqqqff] (3.,3.)-- (-1.,5.);
\draw (-1.52,4.52) node[anchor=north west] {$5$};
\draw [line width=1.pt,color=eqeqeq] (-2.,2.)-- (4.,2.);
\draw [line width=1.pt,color=eqeqeq] (4.,2.)-- (4.,6.);
\draw [line width=1.pt,color=eqeqeq] (4.,6.)-- (-2.,6.);
\draw [line width=1.pt,color=eqeqeq] (-2.,6.)-- (-2.,2.);
\begin{scriptsize}
\draw [color=ududff] (-1.,5.)-- ++(-2.5pt,0 pt) -- ++(5.0pt,0 pt) ++(-2.5pt,-2.5pt) -- ++(0 pt,5.0pt);
\draw[color=ududff] (-1.1,5.39) node {$A$};
\draw [color=ududff] (-1.,3.)-- ++(-2.5pt,0 pt) -- ++(5.0pt,0 pt) ++(-2.5pt,-2.5pt) -- ++(0 pt,5.0pt);
\draw[color=ududff] (-1.28,2.97) node {$B$};
\draw [color=ududff] (3.,3.)-- ++(-2.5pt,0 pt) -- ++(5.0pt,0 pt) ++(-2.5pt,-2.5pt) -- ++(0 pt,5.0pt);
\draw[color=ududff] (3.14,3.37) node {$C$};
\end{scriptsize}
\end{tikzpicture}
\end{minipage}

\ExoCd{Calculer.}

\begin{minipage}{0.7\linewidth}

Calcule l'aire du triangle ABC tracé ci-contre. L'unité est le centimètre.

\point{3}

Ta réponse : $ {A}_{ABC}=\ldots\ldots ~ \text{cm}^2$

\end{minipage}
\begin{minipage}{0.3\linewidth}
\definecolor{qqwuqq}{rgb}{0.,0.39215686274509803,0.}
\definecolor{eqeqeq}{rgb}{0.8784313725490196,0.8784313725490196,0.8784313725490196}
\definecolor{qqqqff}{rgb}{0.,0.,1.}
\definecolor{ududff}{rgb}{0.30196078431372547,0.30196078431372547,1.}
\begin{tikzpicture}[line cap=round,line join=round,>=triangle 45,x=0.6309148264984229cm,y=0.6756756756756758cm]
\clip(-1.42,1.72) rectangle (5.22,6.34);
\fill[line width=1.pt,color=eqeqeq,fill=eqeqeq,fill opacity=0.10000000149011612] (-1.2,2.) -- (4.8,2.) -- (4.8,6.) -- (-1.2,6.) -- cycle;
\draw[line width=1.pt,color=qqwuqq,fill=qqwuqq,fill opacity=0.10000000149011612] (1.5042640687119284,3.) -- (1.5042640687119284,3.424264068711928) -- (1.08,3.424264068711928) -- (1.08,3.) -- cycle; 
\draw [line width=1.pt,color=qqqqff] (1.08,5.22)-- (-0.2,3.);
\draw [line width=1.pt,color=qqqqff] (-0.2,3.)-- (3.8,3.);
\draw [line width=1.pt,color=qqqqff] (3.8,3.)-- (1.08,5.22);
\draw (1.42,3.16) node[anchor=north west] {$8$};
\draw (1.18,4.36) node[anchor=north west] {$5$};
\draw [line width=1.pt,color=eqeqeq] (-1.2,2.)-- (4.8,2.);
\draw [line width=1.pt,color=eqeqeq] (4.8,2.)-- (4.8,6.);
\draw [line width=1.pt,color=eqeqeq] (4.8,6.)-- (-1.2,6.);
\draw [line width=1.pt,color=eqeqeq] (-1.2,6.)-- (-1.2,2.);
\draw [line width=1.pt,dash pattern=on 2pt off 2pt,color=qqqqff] (1.08,5.22)-- (1.08,3.);
\begin{scriptsize}
\draw [color=ududff] (1.08,5.22)-- ++(-2.5pt,0 pt) -- ++(5.0pt,0 pt) ++(-2.5pt,-2.5pt) -- ++(0 pt,5.0pt);
\draw[color=ududff] (0.98,5.61) node {$A$};
\draw [color=ududff] (-0.2,3.)-- ++(-2.5pt,0 pt) -- ++(5.0pt,0 pt) ++(-2.5pt,-2.5pt) -- ++(0 pt,5.0pt);
\draw[color=ududff] (-0.48,2.97) node {$B$};
\draw [color=ududff] (3.8,3.)-- ++(-2.5pt,0 pt) -- ++(5.0pt,0 pt) ++(-2.5pt,-2.5pt) -- ++(0 pt,5.0pt);
\draw[color=ududff] (3.94,3.37) node {$C$};
\draw [color=qqqqff] (1.08,3.)-- ++(-2.0pt,0 pt) -- ++(4.0pt,0 pt) ++(-2.0pt,-2.0pt) -- ++(0 pt,4.0pt);
\draw[color=qqqqff] (0.88,2.69) node {$J$};
\end{scriptsize}
\end{tikzpicture}

\end{minipage}




\ExoCd{Calculer.}

\begin{minipage}{0.7\linewidth}
Calcule l'aire du triangle ABC tracé ci-contre. L'unité est le centimètre.

\point{3}

Ta réponse : $ {A}_{ABC}=\ldots\ldots ~ \text{cm}^2$

\end{minipage}
\begin{minipage}{0.3\linewidth}


\definecolor{qqttqq}{rgb}{0.,0.2,0.}
\definecolor{uuuuuu}{rgb}{0.26666666666666666,0.26666666666666666,0.26666666666666666}
\definecolor{eqeqeq}{rgb}{0.8784313725490196,0.8784313725490196,0.8784313725490196}
\definecolor{qqqqff}{rgb}{0.,0.,1.}
\definecolor{ududff}{rgb}{0.30196078431372547,0.30196078431372547,1.}
\begin{tikzpicture}[line cap=round,line join=round,>=triangle 45,x=0.6309148264984229cm,y=0.6756756756756758cm]
\clip(-2.5554521975987314,1.5529156395832784) rectangle (4.94531164623033,6.434365125249804);
\fill[line width=1.pt,color=eqeqeq,fill=eqeqeq,fill opacity=0.10000000149011612] (-1.98,1.98) -- (4.02,1.98) -- (4.02,5.98) -- (-1.98,5.98) -- cycle;
\draw[line width=1.2pt,color=qqttqq] (-0.788641601320764,2.9097800098107687) -- (-0.8466321301047103,3.3267061186804527) -- (-1.263558238974394,3.268715589896506) -- (-1.2055677101904476,2.8517894810268225) -- cycle; 
\draw [line width=1.pt,color=qqqqff] (-1.56,5.4)-- (-0.14,3.);
\draw [line width=1.pt,color=qqqqff] (-0.14,3.)-- (3.4372215612276618,3.4975581176130324);
\draw [line width=1.pt,color=qqqqff] (3.4372215612276618,3.4975581176130324)-- (-1.56,5.4);
\draw (1.4727357926057645,3.4578715364287516) node[anchor=north west] {$12$};
\draw (-1.7418772833209761,4.251603160114366) node[anchor=north west] {$7$};
\draw [line width=1.pt,color=eqeqeq] (-1.98,1.98)-- (4.02,1.98);
\draw [line width=1.pt,color=eqeqeq] (4.02,1.98)-- (4.02,5.98);
\draw [line width=1.pt,color=eqeqeq] (4.02,5.98)-- (-1.98,5.98);
\draw [line width=1.pt,color=eqeqeq] (-1.98,5.98)-- (-1.98,1.98);
\draw [line width=1.pt,dash pattern=on 2pt off 2pt,color=qqqqff] (-1.56,5.4)-- (-1.2055677101904476,2.8517894810268225);
\draw [line width=1.pt,dash pattern=on 2pt off 2pt,color=qqqqff] (-1.2055677101904476,2.8517894810268225)-- (-0.14,3.);
\begin{scriptsize}
\draw [color=ududff] (-1.56,5.4)-- ++(-2.5pt,0 pt) -- ++(5.0pt,0 pt) ++(-2.5pt,-2.5pt) -- ++(0 pt,5.0pt);
\draw[color=ududff] (-1.6625041209524145,5.789458181005242) node {$A$};
\draw [color=ududff] (-0.14,3.)-- ++(-2.5pt,0 pt) -- ++(5.0pt,0 pt) ++(-2.5pt,-2.5pt) -- ++(0 pt,5.0pt);
\draw[color=ududff] (-0.015511001804763627,2.7534347204077694) node {$B$};
\draw [color=ududff] (3.4372215612276618,3.4975581176130324)-- ++(-2.5pt,0 pt) -- ++(5.0pt,0 pt) ++(-2.5pt,-2.5pt) -- ++(0 pt,5.0pt);
\draw[color=uuuuuu] (-1.42438463384673,2.495471942709945) node {$L$};
\end{scriptsize}
\end{tikzpicture}

\end{minipage}


\subsection{ Calculer l'aire d'un disque à l'aide de la formule}

\ExoCu{Calculer.}

\begin{minipage}{0.7\linewidth}
Calcule l'aire du carré ABCD tracé ci-contre. L'unité est le centimètre.
\point{3}

Ta réponse : $ {A}_{ABCD}=\ldots\ldots ~ \text{cm}^2$

\end{minipage}
\begin{minipage}{0.3\linewidth}

\definecolor{yqqqyq}{rgb}{0.5019607843137255,0.,0.5019607843137255}
\definecolor{eqeqeq}{rgb}{0.8784313725490196,0.8784313725490196,0.8784313725490196}
\definecolor{xfqqff}{rgb}{0.4980392156862745,0.,1.}
\begin{tikzpicture}[line cap=round,line join=round,>=triangle 45,x=0.6cm,y=0.6cm]
\clip(-2.257802838716626,1.6719753831361204) rectangle (4.508759253203242,6.156559056959839);
\fill[line width=1.pt,color=eqeqeq,fill=eqeqeq,fill opacity=0.10000000149011612] (-1.98,1.98) -- (4.02,1.98) -- (4.02,5.98) -- (-1.98,5.98) -- cycle;
\fill[line width=1.pt,color=yqqqyq,fill=yqqqyq,fill opacity=0.10000000149011612] (-0.6703395913453962,2.267274100900331) -- (2.762549681094888,2.2871173914924716) -- (2.742706390502748,5.720006663932755) -- (-0.6901828819375357,5.700163373340615) -- cycle;
\draw[line width=1.pt,color=xfqqff,fill=xfqqff,fill opacity=0.10000000149011612] (-0.2494068633122391,2.269707238056592) -- (-0.2518400004685001,2.6906399660897495) -- (-0.6727727285016571,2.688206828933488) -- (-0.6703395913453962,2.267274100900331) -- cycle; 
\draw[line width=1.pt,color=xfqqff,fill=xfqqff,fill opacity=0.10000000149011612] (-0.6877497447812748,5.279230645307458) -- (-0.2668170167481177,5.281663782463719) -- (-0.2692501539043786,5.702596510496876) -- (-0.6901828819375357,5.700163373340615) -- cycle; 
\draw[line width=1.pt,color=xfqqff,fill=xfqqff,fill opacity=0.10000000149011612] (2.760116543938627,2.7080501195256286) -- (2.3391838159054696,2.7056169823693677) -- (2.341616953061731,2.2846842543362107) -- (2.762549681094888,2.2871173914924716) -- cycle; 
\draw[line width=1.pt,color=xfqqff,fill=xfqqff,fill opacity=0.10000000149011612] (-0.2494068633122391,2.269707238056592) -- (-0.2518400004685001,2.6906399660897495) -- (-0.6727727285016571,2.688206828933488) -- (-0.6703395913453962,2.267274100900331) -- cycle; 
\draw[line width=1.pt,color=xfqqff,fill=xfqqff,fill opacity=0.10000000149011612] (-0.6877497447812748,5.279230645307458) -- (-0.2668170167481177,5.281663782463719) -- (-0.2692501539043786,5.702596510496876) -- (-0.6901828819375357,5.700163373340615) -- cycle; 
\draw [color=xfqqff](-1.2259517279253265,4.450036066035769) node[anchor=north west] {$9$};
\draw [line width=1.pt,color=eqeqeq] (-1.98,1.98)-- (4.02,1.98);
\draw [line width=1.pt,color=eqeqeq] (4.02,1.98)-- (4.02,5.98);
\draw [line width=1.pt,color=eqeqeq] (4.02,5.98)-- (-1.98,5.98);
\draw [line width=1.pt,color=eqeqeq] (-1.98,5.98)-- (-1.98,1.98);
\draw [line width=1.pt,color=yqqqyq] (-0.6703395913453962,2.267274100900331)-- (2.762549681094888,2.2871173914924716);
\draw [line width=1.pt,color=yqqqyq] (0.995809452025714,2.395966752798848) -- (0.9971858422519138,2.15785124366637);
\draw [line width=1.pt,color=yqqqyq] (1.095024247497579,2.396540248726432) -- (1.0964006377237787,2.1584247395939538);
\draw [line width=1.pt,color=yqqqyq] (2.762549681094888,2.2871173914924716)-- (2.742706390502748,5.720006663932755);
\draw [line width=1.pt,color=yqqqyq] (2.63385702919637,3.95326643486358) -- (2.8719725383288495,3.95464282508978);
\draw [line width=1.pt,color=yqqqyq] (2.633283533268788,4.052481230335446) -- (2.8713990424012668,4.0538576205616454);
\draw [line width=1.pt,color=yqqqyq] (2.742706390502748,5.720006663932755)-- (-0.6901828819375357,5.700163373340615);
\draw [line width=1.pt,color=yqqqyq] (1.0765573471316383,5.591314012034237) -- (1.0751809569054387,5.829429521166716);
\draw [line width=1.pt,color=yqqqyq] (0.9773425516597734,5.590740516106655) -- (0.9759661614335737,5.828856025239133);
\draw [line width=1.pt,color=yqqqyq] (-0.6901828819375357,5.700163373340615)-- (-0.6703395913453962,2.267274100900331);
\draw [line width=1.pt,color=yqqqyq] (-0.5614902300390178,4.034014329969506) -- (-0.7996057391714969,4.0326379397433065);
\draw [line width=1.pt,color=yqqqyq] (-0.5609167341114354,3.93479953449764) -- (-0.7990322432439145,3.93342314427144);
\begin{scriptsize}
\draw [color=xfqqff] (-0.6703395913453962,2.267274100900331)-- ++(-2.5pt,0 pt) -- ++(5.0pt,0 pt) ++(-2.5pt,-2.5pt) -- ++(0 pt,5.0pt);
\draw[color=xfqqff] (-1.1465785655567649,2.3764121991571026) node {$A$};
\draw [color=xfqqff] (2.762549681094888,2.2871173914924716)-- ++(-2.5pt,0 pt) -- ++(5.0pt,0 pt) ++(-2.5pt,-2.5pt) -- ++(0 pt,5.0pt);
\draw[color=xfqqff] (2.9014527152398717,2.654218267447068) node {$B$};
\draw [color=xfqqff] (2.742706390502748,5.720006663932755)-- ++(-2.5pt,0 pt) -- ++(5.0pt,0 pt) ++(-2.5pt,-2.5pt) -- ++(0 pt,5.0pt);
\draw[color=xfqqff] (3.0205124587927137,5.6902417280445405) node {$C$};
\draw [color=xfqqff] (-0.6901828819375357,5.700163373340615)-- ++(-2.5pt,0 pt) -- ++(5.0pt,0 pt) ++(-2.5pt,-2.5pt) -- ++(0 pt,5.0pt);
\draw[color=xfqqff] (-1.1465785655567649,5.749771599820961) node {$D$};
\end{scriptsize}
\end{tikzpicture}

\end{minipage}


\ExoCd{Calculer.}

\begin{minipage}{0.7\linewidth}
L'aire du carré ABCD mesure $49$ cm$^2$. Quelle est la longueur du coté ?
\point{3}

Ta réponse : $ {A}_{ABCD}=\ldots\ldots ~ \text{cm}^2$

\end{minipage}
\begin{minipage}{0.3\linewidth}

\definecolor{yqqqyq}{rgb}{0.5019607843137255,0.,0.5019607843137255}
\definecolor{eqeqeq}{rgb}{0.8784313725490196,0.8784313725490196,0.8784313725490196}
\definecolor{xfqqff}{rgb}{0.4980392156862745,0.,1.}
\begin{tikzpicture}[line cap=round,line join=round,>=triangle 45,x=0.6cm,y=0.6cm]
\clip(-2.257802838716626,1.6719753831361204) rectangle (4.508759253203242,6.156559056959839);
\fill[line width=1.pt,color=eqeqeq,fill=eqeqeq,fill opacity=0.10000000149011612] (-1.98,1.98) -- (4.02,1.98) -- (4.02,5.98) -- (-1.98,5.98) -- cycle;
\fill[line width=1.pt,color=yqqqyq,fill=yqqqyq,fill opacity=0.10000000149011612] (-0.6703395913453962,2.267274100900331) -- (2.762549681094888,2.2871173914924716) -- (2.742706390502748,5.720006663932755) -- (-0.6901828819375357,5.700163373340615) -- cycle;
\draw[line width=1.pt,color=xfqqff,fill=xfqqff,fill opacity=0.10000000149011612] (-0.2494068633122391,2.269707238056592) -- (-0.2518400004685001,2.6906399660897495) -- (-0.6727727285016571,2.688206828933488) -- (-0.6703395913453962,2.267274100900331) -- cycle; 
\draw[line width=1.pt,color=xfqqff,fill=xfqqff,fill opacity=0.10000000149011612] (-0.6877497447812748,5.279230645307458) -- (-0.2668170167481177,5.281663782463719) -- (-0.2692501539043786,5.702596510496876) -- (-0.6901828819375357,5.700163373340615) -- cycle; 
\draw[line width=1.pt,color=xfqqff,fill=xfqqff,fill opacity=0.10000000149011612] (2.760116543938627,2.7080501195256286) -- (2.3391838159054696,2.7056169823693677) -- (2.341616953061731,2.2846842543362107) -- (2.762549681094888,2.2871173914924716) -- cycle; 
\draw[line width=1.pt,color=xfqqff,fill=xfqqff,fill opacity=0.10000000149011612] (-0.2494068633122391,2.269707238056592) -- (-0.2518400004685001,2.6906399660897495) -- (-0.6727727285016571,2.688206828933488) -- (-0.6703395913453962,2.267274100900331) -- cycle; 
\draw[line width=1.pt,color=xfqqff,fill=xfqqff,fill opacity=0.10000000149011612] (-0.6877497447812748,5.279230645307458) -- (-0.2668170167481177,5.281663782463719) -- (-0.2692501539043786,5.702596510496876) -- (-0.6901828819375357,5.700163373340615) -- cycle; 
\draw [line width=1.pt,color=eqeqeq] (-1.98,1.98)-- (4.02,1.98);
\draw [line width=1.pt,color=eqeqeq] (4.02,1.98)-- (4.02,5.98);
\draw [line width=1.pt,color=eqeqeq] (4.02,5.98)-- (-1.98,5.98);
\draw [line width=1.pt,color=eqeqeq] (-1.98,5.98)-- (-1.98,1.98);
\draw [line width=1.pt,color=yqqqyq] (-0.6703395913453962,2.267274100900331)-- (2.762549681094888,2.2871173914924716);
\draw [line width=1.pt,color=yqqqyq] (0.995809452025714,2.395966752798848) -- (0.9971858422519138,2.15785124366637);
\draw [line width=1.pt,color=yqqqyq] (1.095024247497579,2.396540248726432) -- (1.0964006377237787,2.1584247395939538);
\draw [line width=1.pt,color=yqqqyq] (2.762549681094888,2.2871173914924716)-- (2.742706390502748,5.720006663932755);
\draw [line width=1.pt,color=yqqqyq] (2.63385702919637,3.95326643486358) -- (2.8719725383288495,3.95464282508978);
\draw [line width=1.pt,color=yqqqyq] (2.633283533268788,4.052481230335446) -- (2.8713990424012668,4.0538576205616454);
\draw [line width=1.pt,color=yqqqyq] (2.742706390502748,5.720006663932755)-- (-0.6901828819375357,5.700163373340615);
\draw [line width=1.pt,color=yqqqyq] (1.0765573471316383,5.591314012034237) -- (1.0751809569054387,5.829429521166716);
\draw [line width=1.pt,color=yqqqyq] (0.9773425516597734,5.590740516106655) -- (0.9759661614335737,5.828856025239133);
\draw [line width=1.pt,color=yqqqyq] (-0.6901828819375357,5.700163373340615)-- (-0.6703395913453962,2.267274100900331);
\draw [line width=1.pt,color=yqqqyq] (-0.5614902300390178,4.034014329969506) -- (-0.7996057391714969,4.0326379397433065);
\draw [line width=1.pt,color=yqqqyq] (-0.5609167341114354,3.93479953449764) -- (-0.7990322432439145,3.93342314427144);
\begin{scriptsize}
\draw [color=xfqqff] (-0.6703395913453962,2.267274100900331)-- ++(-2.5pt,0 pt) -- ++(5.0pt,0 pt) ++(-2.5pt,-2.5pt) -- ++(0 pt,5.0pt);
\draw[color=xfqqff] (-1.1465785655567649,2.3764121991571026) node {$A$};
\draw [color=xfqqff] (2.762549681094888,2.2871173914924716)-- ++(-2.5pt,0 pt) -- ++(5.0pt,0 pt) ++(-2.5pt,-2.5pt) -- ++(0 pt,5.0pt);
\draw[color=xfqqff] (2.9014527152398717,2.654218267447068) node {$B$};
\draw [color=xfqqff] (2.742706390502748,5.720006663932755)-- ++(-2.5pt,0 pt) -- ++(5.0pt,0 pt) ++(-2.5pt,-2.5pt) -- ++(0 pt,5.0pt);
\draw[color=xfqqff] (3.0205124587927137,5.6902417280445405) node {$C$};
\draw [color=xfqqff] (-0.6901828819375357,5.700163373340615)-- ++(-2.5pt,0 pt) -- ++(5.0pt,0 pt) ++(-2.5pt,-2.5pt) -- ++(0 pt,5.0pt);
\draw[color=xfqqff] (-1.1465785655567649,5.749771599820961) node {$D$};
\end{scriptsize}
\end{tikzpicture}

\end{minipage}




\ExoCu{Calculer.}

\begin{minipage}{0.7\linewidth}
Calcule l'aire du disque de centre A et de rayon $5$ cm tracé ci-contre. 

\point{3}

Ta réponse : $ {A}=\ldots\ldots ~ \text{cm}^2$

\end{minipage}
\begin{minipage}{0.3\linewidth}

 
\definecolor{ffqqff}{rgb}{1.,0.,1.}
\definecolor{ududff}{rgb}{0.30196078431372547,0.30196078431372547,1.}
\definecolor{eqeqeq}{rgb}{0.8784313725490196,0.8784313725490196,0.8784313725490196}
\definecolor{xfqqff}{rgb}{0.4980392156862745,0.,1.}
\begin{tikzpicture}[line cap=round,line join=round,>=triangle 45,x=0.5999388987165869cm,y=0.5999388987165869cm]
\clip(-2.218116257532345,1.5529156395832784) rectangle (4.449229381426821,6.295462091104821);
\fill[line width=1.pt,color=eqeqeq,fill=eqeqeq,fill opacity=0.10000000149011612] (-1.99984329059214,1.860940256447158) -- (4.000156709407859,1.860940256447158) -- (4.000156709407859,5.860940256447158) -- (-1.99984329059214,5.860940256447158) -- cycle;
\draw [color=xfqqff](1.4727357926057645,3.715834314126576) node[anchor=north west] {$5$};
\draw [line width=1.pt,color=eqeqeq] (-1.99984329059214,1.860940256447158)-- (4.000156709407859,1.860940256447158);
\draw [line width=1.pt,color=eqeqeq] (4.000156709407859,1.860940256447158)-- (4.000156709407859,5.860940256447158);
\draw [line width=1.pt,color=eqeqeq] (4.000156709407859,5.860940256447158)-- (-1.99984329059214,5.860940256447158);
\draw [line width=1.pt,color=eqeqeq] (-1.99984329059214,5.860940256447158)-- (-1.99984329059214,1.860940256447158);
\draw [line width=1.pt,color=ffqqff] (1.1552431431315184,3.9142672200479796) circle (1.1036170176025513cm);
\draw [line width=1.pt,color=xfqqff] (1.1552431431315184,3.9142672200479796)-- (2.8419228434634505,3.1800654681387868);
\begin{scriptsize}
\draw [color=ududff] (1.1552431431315184,3.9142672200479796)-- ++(-2.0pt,0 pt) -- ++(4.0pt,0 pt) ++(-2.0pt,-2.0pt) -- ++(0 pt,4.0pt);
\draw[color=ududff] (1.1155565619472378,4.360741258371137) node {$A$};
\draw [color=ududff] (2.8419228434634505,3.1800654681387868)-- ++(-0.5pt,0 pt) -- ++(1.0pt,0 pt) ++(-0.5pt,-0.5pt) -- ++(0 pt,1.0pt);
\draw[color=ududff] (3.0205124587927137,3.249516985211278) node {$B$};
\end{scriptsize}
\end{tikzpicture}

\end{minipage}




\ExoCu{Calculer.}

\begin{minipage}{0.7\linewidth}
Calcule l'aire du disque de centre A et de diamètre $19$ cm tracé ci-contre. 

\point{3}

Ta réponse : $ {A}=\ldots\ldots ~ \text{cm}^2$

\end{minipage}
\begin{minipage}{0.3\linewidth}

 

\definecolor{ubqqys}{rgb}{0.29411764705882354,0.,0.5098039215686274}
\definecolor{xdxdff}{rgb}{0.49019607843137253,0.49019607843137253,1.}
\definecolor{qqttcc}{rgb}{0.,0.2,0.8}
\definecolor{ududff}{rgb}{0.30196078431372547,0.30196078431372547,1.}
\definecolor{eqeqeq}{rgb}{0.8784313725490196,0.8784313725490196,0.8784313725490196}
\definecolor{xfqqff}{rgb}{0.4980392156862745,0.,1.}
\begin{tikzpicture}[line cap=round,line join=round,>=triangle 45,x=0.6338977043043182cm,y=0.6338977043043182cm]
\clip(-2.059369932795222,1.7315052549125414) rectangle (4.250796475505417,6.176402347551979);
\fill[line width=1.pt,color=eqeqeq,fill=eqeqeq,fill opacity=0.10000000149011612] (-1.9403134188157194,1.9403134188157196) -- (4.05968658118428,1.9403134188157196) -- (4.05968658118428,5.94031341881572) -- (-1.9403134188157194,5.94031341881572) -- cycle;
\draw [color=xfqqff](0.5004145535908862,3.894423929455839) node[anchor=north west] {$19$};
\draw [line width=1.pt,color=eqeqeq] (-1.9403134188157194,1.9403134188157196)-- (4.05968658118428,1.9403134188157196);
\draw [line width=1.pt,color=eqeqeq] (4.05968658118428,1.9403134188157196)-- (4.05968658118428,5.94031341881572);
\draw [line width=1.pt,color=eqeqeq] (4.05968658118428,5.94031341881572)-- (-1.9403134188157194,5.94031341881572);
\draw [line width=1.pt,color=eqeqeq] (-1.9403134188157194,5.94031341881572)-- (-1.9403134188157194,1.9403134188157196);
\draw [line width=1.pt,color=qqttcc] (1.1552431431315184,3.9142672200479796) circle (1.166085905391375cm);
\draw [line width=1.pt,color=ubqqys] (2.8419228434634505,3.1800654681387868)-- (-0.539765658736616,4.629029967823699);
\begin{scriptsize}
\draw [color=ududff] (1.1552431431315184,3.9142672200479796)-- ++(-2.0pt,0 pt) -- ++(4.0pt,0 pt) ++(-2.0pt,-2.0pt) -- ++(0 pt,4.0pt);
\draw[color=ududff] (1.2147730149079397,4.301211386594716) node {$A$};
\draw [color=ududff] (2.8419228434634505,3.1800654681387868)-- ++(-1.5pt,0 pt) -- ++(3.0pt,0 pt) ++(-1.5pt,-1.5pt) -- ++(0 pt,3.0pt);
\draw[color=ududff] (3.0205124587927137,3.3288901475798394) node {$B$};
\draw [color=xdxdff] (-0.539765658736616,4.629029967823699)-- ++(-2.5pt,0 pt) -- ++(5.0pt,0 pt) ++(-2.5pt,-2.5pt) -- ++(0 pt,5.0pt);
\draw[color=xdxdff] (-1.0275188220039226,4.698077198437523) node {$C$};
\end{scriptsize}
\end{tikzpicture}
\end{minipage}

\ExoCd{Calculer.}

\begin{minipage}{0.7\linewidth}

On a tracé un cercle de centre A et deux cercles de centre D et E dont les rayons mesures $6$ cm.
Calcule l'aire de la partie hachurée. 

\point{5}

Ta réponse : $ {A}=\ldots\ldots ~ \text{cm}^2$

\end{minipage}
\begin{minipage}{0.3\linewidth}

 
 
\definecolor{qqqqff}{rgb}{0.,0.,1.}
\definecolor{xfqqff}{rgb}{0.4980392156862745,0.,1.}
\definecolor{ffffff}{rgb}{1.,1.,1.}
\definecolor{uuuuuu}{rgb}{0.26666666666666666,0.26666666666666666,0.26666666666666666}
\definecolor{xdxdff}{rgb}{0.49019607843137253,0.49019607843137253,1.}
\definecolor{qqttcc}{rgb}{0.,0.2,0.8}
\definecolor{ududff}{rgb}{0.30196078431372547,0.30196078431372547,1.}
\definecolor{eqeqeq}{rgb}{0.8784313725490196,0.8784313725490196,0.8784313725490196}
\begin{tikzpicture}[line cap=round,line join=round,>=triangle 45,x=1.0cm,y=1.0cm]
\clip(-1.5236010868074321,1.7513485455046818) rectangle (4.012676988399733,6.097029185183418);
\fill[line width=1.pt,color=eqeqeq,fill=eqeqeq,fill opacity=0.10000000149011612] (-1.2976493588821056,1.9007835470392984) -- (3.7023506411178944,1.9007835470392984) -- (3.7023506411178944,5.900783547039298) -- (-1.2976493588821056,5.900783547039298) -- cycle;
\draw [line width=1.pt,color=qqttcc,fill=qqttcc,pattern=north east lines,pattern color=qqttcc] (1.1552431431315184,3.9142672200479796) circle (1.8395490273483728cm);
\draw [line width=1.pt,color=ffffff,fill=ffffff,fill opacity=1.0] (1.9985829932974846,3.547166344093383) circle (0.9197745136741862cm);
\draw [line width=1.pt,color=ffffff,fill=ffffff,fill opacity=1.0] (0.30773874219745123,4.271648593935839) circle (0.9197745136741864cm);
\draw [line width=1.pt,color=eqeqeq] (-1.2976493588821056,1.9007835470392984)-- (3.7023506411178944,1.9007835470392984);
\draw [line width=1.pt,color=eqeqeq] (3.7023506411178944,1.9007835470392984)-- (3.7023506411178944,5.900783547039298);
\draw [line width=1.pt,color=eqeqeq] (3.7023506411178944,5.900783547039298)-- (-1.2976493588821056,5.900783547039298);
\draw [line width=1.pt,color=eqeqeq] (-1.2976493588821056,5.900783547039298)-- (-1.2976493588821056,1.9007835470392984);
\draw [line width=1.pt,color=xfqqff] (-0.539765658736616,4.629029967823699)-- (0.30773874219745123,4.271648593935839);
\draw [line width=1.pt,color=xfqqff] (-0.06975240831324084,4.56004405649052) -- (-0.16227450822592343,4.3406345052690165);
\draw [line width=1.pt,color=xfqqff] (0.30773874219745123,4.271648593935839)-- (1.1552431431315184,3.9142672200479796);
\draw [line width=1.pt,color=xfqqff] (0.7777519926208263,4.202662682602661) -- (0.6852298927081436,3.983253131381157);
\draw [line width=1.pt,color=xfqqff] (1.1552431431315184,3.9142672200479796)-- (1.9985829932974846,3.547166344093383);
\draw [line width=1.pt,color=xfqqff] (1.6244322541664842,3.8398824795279385) -- (1.5293938822625186,3.6215510846134245);
\draw [line width=1.pt,color=xfqqff] (1.9985829932974846,3.547166344093383)-- (2.8419228434634505,3.1800654681387868);
\draw [line width=1.pt,color=xfqqff] (2.46777210433245,3.472781603573342) -- (2.3727337324284847,3.254450208658828);
\draw [line width=1.pt,color=qqqqff] (1.3378001612497492,2.9073607448090693)-- (1.9985829932974846,3.547166344093383);
\draw [color=xfqqff](1.611638826750747,3.4975581176130324) node[anchor=north west] {$6 ~\text{cm}$};
\begin{scriptsize}
\draw [color=ududff] (1.1552431431315184,3.9142672200479796)-- ++(-2.0pt,0 pt) -- ++(4.0pt,0 pt) ++(-2.0pt,-2.0pt) -- ++(0 pt,4.0pt);
\draw[color=ududff] (1.4925790831979049,4.142465061857593) node {$A$};
\draw [color=ududff] (2.8419228434634505,3.1800654681387868)-- ++(-1.5pt,0 pt) -- ++(3.0pt,0 pt) ++(-1.5pt,-1.5pt) -- ++(0 pt,3.0pt);
\draw[color=ududff] (3.0205124587927137,3.3288901475798394) node {$B$};
\draw [color=xdxdff] (-0.539765658736616,4.629029967823699)-- ++(-2.5pt,0 pt) -- ++(5.0pt,0 pt) ++(-2.5pt,-2.5pt) -- ++(0 pt,5.0pt);
\draw[color=xdxdff] (-0.8489292066746593,4.817136941990365) node {$C$};
\draw [color=uuuuuu] (0.30773874219745123,4.271648593935839)-- ++(-2.0pt,0 pt) -- ++(4.0pt,0 pt) ++(-2.0pt,-2.0pt) -- ++(0 pt,4.0pt);
\draw[color=uuuuuu] (0.440884681814465,4.598860745476821) node {$D$};
\draw [color=uuuuuu] (1.9985829932974846,3.547166344093383)-- ++(-2.0pt,0 pt) -- ++(4.0pt,0 pt) ++(-2.0pt,-2.0pt) -- ++(0 pt,4.0pt);
\draw[color=uuuuuu] (2.1275643821463968,3.983718737120471) node {$E$};
\end{scriptsize}
\end{tikzpicture}

\end{minipage}




\ExoCd{Calculer.}


Calcule l'aire d'un disque de périmètre $5 \pi$ cm. 

\point{3}

Ta réponse : $ {A}_{ABC}=\ldots\ldots ~ \text{cm}^2$

 
 
 
 
 
 
 

\subsection{ Déterminer la mesure de l’aire d’une surface}



\ExoCu{Calculer.}

\begin{minipage}{0.7\linewidth}

Détermine l'aire de la surface violette.

\point{5}

Ta réponse : $\mathcal{A}_{violette}=\ldots\ldots ~ \text{cm}^2$

\end{minipage}
\begin{minipage}{0.3\linewidth}

\definecolor{ffffff}{rgb}{1.,1.,1.}
\definecolor{uuuuuu}{rgb}{0.26666666666666666,0.26666666666666666,0.26666666666666666}
\definecolor{xfqqff}{rgb}{0.4980392156862745,0.,1.}
\definecolor{eqeqeq}{rgb}{0.8784313725490196,0.8784313725490196,0.8784313725490196}
\definecolor{qqqqff}{rgb}{0.,0.,1.}
\begin{tikzpicture}[line cap=round,line join=round,>=triangle 45,x=0.8726383981332172cm,y=0.8726383981332172cm]
\clip(-2.317332710493047,1.7910351266889624) rectangle (2.2664674162913796,6.25577550992054);
\fill[line width=1.pt,color=eqeqeq,fill=eqeqeq,fill opacity=0.15000000596046448] (-2.,2.) -- (2.,2.) -- (2.,6.) -- (-2.,6.) -- cycle;
\fill[line width=1.pt,color=xfqqff,fill=xfqqff,fill opacity=0.10000000149011612] (-1.4045413432545897,2.5450801691902956) -- (1.532265664382186,2.5450801691902956) -- (1.5322656643821861,5.481887176827071) -- (-1.4045413432545895,5.481887176827072) -- cycle;
\draw [line width=1.pt,color=ffffff,fill=ffffff,fill opacity=1.0] (0.06386216056379823,4.013483673008683) circle (1.2813852813852815cm);
\draw [color=qqqqff](-0.9283023690432208,4.072229997745804) node[anchor=north west] {$10$ cm};
\draw [line width=1.pt,color=eqeqeq] (-2.,2.)-- (2.,2.);
\draw [line width=1.pt,color=eqeqeq] (2.,2.)-- (2.,6.);
\draw [line width=1.pt,color=eqeqeq] (2.,6.)-- (-2.,6.);
\draw [line width=1.pt,color=eqeqeq] (-2.,6.)-- (-2.,2.);
\draw [line width=1.pt,color=xfqqff] (-1.4045413432545897,2.5450801691902956)-- (1.532265664382186,2.5450801691902956);
\draw [line width=1.pt,color=xfqqff] (1.532265664382186,2.5450801691902956)-- (1.5322656643821861,5.481887176827071);
\draw [line width=1.pt,color=xfqqff] (1.5322656643821861,5.481887176827071)-- (-1.4045413432545895,5.481887176827072);
\draw [line width=1.pt,color=xfqqff] (-1.4045413432545895,5.481887176827072)-- (-1.4045413432545897,2.5450801691902956);
\draw [line width=1.pt,color=qqqqff] (-1.4045413432545897,4.013483673008683)-- (0.06386216056379823,4.013483673008683);
\begin{scriptsize}
\draw [color=xfqqff] (-1.4045413432545897,2.5450801691902956)-- ++(-2.5pt,0 pt) -- ++(5.0pt,0 pt) ++(-2.5pt,-2.5pt) -- ++(0 pt,5.0pt);
\draw[color=xfqqff] (-1.583130958583853,2.3764121991571026) node {$A$};
\draw [color=xfqqff] (1.532265664382186,2.5450801691902956)-- ++(-2.5pt,0 pt) -- ++(5.0pt,0 pt) ++(-2.5pt,-2.5pt) -- ++(0 pt,5.0pt);
\draw[color=xfqqff] (1.6314821173428875,2.316882327380682) node {$B$};
\draw [color=xfqqff] (1.5322656643821861,5.481887176827071)-- ++(-2.5pt,0 pt) -- ++(5.0pt,0 pt) ++(-2.5pt,-2.5pt) -- ++(0 pt,5.0pt);
\draw[color=xfqqff] (1.6711686985271683,5.769614890413101) node {$C$};
\draw [color=xfqqff] (-1.4045413432545895,5.481887176827072)-- ++(-2.5pt,0 pt) -- ++(5.0pt,0 pt) ++(-2.5pt,-2.5pt) -- ++(0 pt,5.0pt);
\draw[color=xfqqff] (-1.6625041209524145,5.6703984374524) node {$D$};
\draw [color=uuuuuu] (0.06386216056379823,4.013483673008683)-- ++(-2.0pt,0 pt) -- ++(4.0pt,0 pt) ++(-2.0pt,-2.0pt) -- ++(0 pt,4.0pt);
\draw[color=uuuuuu] (0.14323532293235935,4.2416815148182945) node {$O$};
\draw[color=ffffff] (-0.6108097195689749,5.035413138503909) node {$c$};
\end{scriptsize}
\end{tikzpicture}
 

\end{minipage}


\ExoCd{Calculer.}

\begin{minipage}{0.7\linewidth}

Calcule la valeur exacte de l'aire de la surface coloré en bleu. 

\point{5}

Ta réponse : $\mathcal{A}_{bleue}=\ldots\ldots ~ \text{cm}^2$

\end{minipage}
\begin{minipage}{0.3\linewidth}


\definecolor{ffffff}{rgb}{1.,1.,1.}
\definecolor{ududff}{rgb}{0.30196078431372547,0.30196078431372547,1.}
\definecolor{eqeqeq}{rgb}{0.8784313725490196,0.8784313725490196,0.8784313725490196}
\definecolor{qqqqff}{rgb}{0.,0.,1.}
\begin{tikzpicture}[line cap=round,line join=round,>=triangle 45,x=0.7999388987165869cm,y=0.7999388987165869cm]
\clip(-2.3768625822694682,1.6124455113596994) rectangle (4.290483056689697,5.243767689721383);
\fill[line width=1.pt,color=eqeqeq,fill=eqeqeq,fill opacity=0.15000000596046448] (-1.99984329059214,1.860940256447158) -- (4.000156709407859,1.860940256447158) -- (4.,5.) -- (-2.,5.) -- cycle;

\draw [line width=1.pt,color=qqqqff] (-1.0076755314117822,2.5847667503745764)-- (2.9609825870162925,2.624453331558857);
\draw [shift={(0.9766535278022551,2.6046100409667168)},line width=1.pt,color=ududff,fill=ududff,pattern=north east lines,pattern color=ududff]  plot[domain=0.009999666686665253:3.151592320276458,variable=\t]({1.*1.9844282731867107*cos(\t r)+0.*1.9844282731867107*sin(\t r)},{0.*1.9844282731867107*cos(\t r)+1.*1.9844282731867107*sin(\t r)});
\draw [line width=1.pt,color=qqqqff] (0.9964968183943955,2.6046100409667168)-- (0.9971104641719741,4.588932868829451);
\draw [shift={(0.9968036412831849,3.596771454898084)},line width=1.pt,color=ffffff,fill=ffffff,fill opacity=1.0]  plot[domain=-1.571105573732643:1.57048707985715,variable=\t]({1.*0.9921614613733869*cos(\t r)+0.*0.9921614613733869*sin(\t r)},{0.*0.9921614613733869*cos(\t r)+1.*0.9921614613733869*sin(\t r)});
\draw [shift={(0.9968036412831849,3.596771454898084)},line width=1.pt,color=ffffff,fill=ffffff,fill opacity=1.0]  plot[domain=1.5704870798571502:4.712079733446943,variable=\t]({1.*0.9921614613733865*cos(\t r)+0.*0.9921614613733865*sin(\t r)},{0.*0.9921614613733865*cos(\t r)+1.*0.9921614613733865*sin(\t r)});
\draw [line width=1.pt,color=eqeqeq] (-1.99984329059214,1.860940256447158)-- (4.000156709407859,1.860940256447158);
\draw [line width=1.pt,color=eqeqeq] (4.000156709407859,1.860940256447158)-- (4.,5.);
\draw [line width=1.pt,color=eqeqeq] (4.,5.)-- (-2.,5.);
\draw [line width=1.pt,color=eqeqeq] (-2.,5.)-- (-1.99984329059214,1.860940256447158);
\draw [color=qqqqff](1.0163401089865358,4.033326963600822) node[anchor=north west] {$10$ cm};
\end{tikzpicture}
 

\end{minipage}






\ExoCd{Calculer.}

\begin{minipage}{0.7\linewidth}

Calcule la valeur exacte de l'aire de la surface violette. 

\point{5}

Ta réponse : $\mathcal{A}_{violette}=\ldots\ldots ~ \text{cm}^2$

\end{minipage}
\begin{minipage}{0.3\linewidth}


\definecolor{eqeqeq}{rgb}{0.8784313725490196,0.8784313725490196,0.8784313725490196}
\definecolor{wwqqcc}{rgb}{0.4,0.,0.8}
\definecolor{aqaqaq}{rgb}{0.6274509803921569,0.6274509803921569,0.6274509803921569}
\definecolor{ffffff}{rgb}{1.,1.,1.}
\definecolor{cqcqcq}{rgb}{0.7529411764705882,0.7529411764705882,0.7529411764705882}
\definecolor{xfqqff}{rgb}{0.4980392156862745,0.,1.}
\definecolor{ududff}{rgb}{0.30196078431372547,0.30196078431372547,1.}
\begin{tikzpicture}[line cap=round,line join=round,>=triangle 45,x=0.743494423791822cm,y=0.743494423791822cm]
\clip(-1.76,1.66) rectangle (3.62,6.8);
\fill[line width=1.pt,color=eqeqeq,fill=eqeqeq,fill opacity=0.10000000149011612] (-1.4,1.78) -- (3.48,1.78) -- (3.48,6.66) -- (-1.4,6.66) -- cycle;
\draw [line width=1.4pt,color=cqcqcq] (1.,4.) circle (1.5339787330032777cm);
\draw [line width=1.pt,color=ffffff] (1.,4.)-- (2.98,3.42);
\draw [line width=1.pt,color=ffffff] (1.9757503004512433,3.8392166632550544) -- (1.9082823352408016,3.6088949889159623);
\draw [line width=1.pt,color=ffffff] (2.071717664759199,3.8111050110840377) -- (2.004249699548757,3.5807833367449455);
\draw (-0.46,4.18) node[anchor=north west] {$8$ cm};
\draw [shift={(1.99,3.71)},line width=1.pt,color=xfqqff,fill=xfqqff,fill opacity=1.0]  plot[domain=2.856635303339639:5.998227956929432,variable=\t]({1.*1.0316006979447037*cos(\t r)+0.*1.0316006979447037*sin(\t r)},{0.*1.0316006979447037*cos(\t r)+1.*1.0316006979447037*sin(\t r)});
\draw [shift={(0.01,4.29)},line width=1.pt,color=ffffff,fill=ffffff,fill opacity=1.0]  plot[domain=-0.28495735025015456:2.856635303339639,variable=\t]({1.*1.0316006979447037*cos(\t r)+0.*1.0316006979447037*sin(\t r)},{0.*1.0316006979447037*cos(\t r)+1.*1.0316006979447037*sin(\t r)});
\draw [line width=1.pt,color=aqaqaq] (-0.98,4.58)-- (1.,4.);
\draw [line width=1.pt,color=aqaqaq] (-0.004249699548757347,4.419216663255054) -- (-0.07171766475919807,4.188894988915963);
\draw [line width=1.pt,color=aqaqaq] (0.09171766475919806,4.391105011084037) -- (0.024249699548757338,4.160783336744945);
\draw [shift={(1.,4.)},line width=1.pt,color=wwqqcc]  plot[domain=-0.28495735025015456:2.856635303339639,variable=\t]({1.*2.0632013958894073*cos(\t r)+0.*2.0632013958894073*sin(\t r)},{0.*2.0632013958894073*cos(\t r)+1.*2.0632013958894073*sin(\t r)});
\draw [line width=1.pt,color=eqeqeq] (-1.4,1.78)-- (3.48,1.78);
\draw [line width=1.pt,color=eqeqeq] (3.48,1.78)-- (3.48,6.66);
\draw [line width=1.pt,color=eqeqeq] (3.48,6.66)-- (-1.4,6.66);
\draw [line width=1.pt,color=eqeqeq] (-1.4,6.66)-- (-1.4,1.78);
\draw [shift={(1.,4.)},line width=1.pt,color=xfqqff,fill=xfqqff,fill opacity=1.0]  plot[domain=-0.28495735025015456:2.856635303339639,variable=\t]({1.*2.0632013958894073*cos(\t r)+0.*2.0632013958894073*sin(\t r)},{0.*2.0632013958894073*cos(\t r)+1.*2.0632013958894073*sin(\t r)});
\end{tikzpicture}

\end{minipage}


\subsection{ Calculer le volume d'un cube ou d'un pavé droit en utilisant une formule}



\ExoCu{Calculer.}

\begin{minipage}{0.7\linewidth}

ABCDEFGH est un cube de coté 8 cm.  

Calcule le volume du cube ABCDEFGH en cm$^3$ et en dm$^3$.

\point{4}

\end{minipage}
\hfill
\begin{minipage}{0.3\linewidth}


\definecolor{ududff}{rgb}{0.30196078431372547,0.30196078431372547,1.}
\definecolor{qqqqff}{rgb}{0.,0.,1.}
\definecolor{eqeqeq}{rgb}{0.8784313725490196,0.8784313725490196,0.8784313725490196}
\definecolor{qqttcc}{rgb}{0.,0.2,0.8}
\begin{tikzpicture}[line cap=round,line join=round,>=triangle 45,x=0.7462686567164183cm,y=0.7462686567164183cm]
\clip(-2.06,1.64) rectangle (3.58,7.1);
\fill[line width=1.pt,color=eqeqeq,fill=eqeqeq,fill opacity=0.10000000149011612] (-1.76,1.86) -- (3.12,1.86) -- (3.12,6.74) -- (-1.76,6.74) -- cycle;
\draw [color=qqttcc](-0.12,3.24) node[anchor=north west] {$8$};
\draw [line width=1.pt,color=eqeqeq] (-1.76,1.86)-- (3.12,1.86);
\draw [line width=1.pt,color=eqeqeq] (3.12,1.86)-- (3.12,6.74);
\draw [line width=1.pt,color=eqeqeq] (3.12,6.74)-- (-1.76,6.74);
\draw [line width=1.pt,color=eqeqeq] (-1.76,6.74)-- (-1.76,1.86);
\draw [line width=1.pt,color=eqeqeq] (-1.,3.)-- (1.,3.);
\draw [line width=1.pt,color=qqqqff] (-1.,3.)-- (1.,3.);
\draw [line width=1.pt,color=qqqqff] (1.,3.)-- (1.,5.);
\draw [line width=1.pt,color=qqqqff] (1.,5.)-- (-1.,5.);
\draw [line width=1.pt,color=qqqqff] (-1.,5.)-- (-1.,3.);
\draw [line width=1.pt,color=qqqqff] (-1.,5.)-- (0.32,5.64);
\draw [line width=1.pt,color=qqqqff] (0.32,5.64)-- (2.32,5.64);
\draw [line width=1.pt,color=qqqqff] (2.32,5.64)-- (2.32,3.64);
\draw [line width=1.pt,color=qqqqff] (2.32,3.64)-- (1.,3.);
\draw [line width=1.pt,color=qqqqff] (1.,5.)-- (2.32,5.64);
\draw [line width=1.pt,dash pattern=on 2pt off 2pt,color=qqqqff] (-1.,3.)-- (0.32,3.64);
\draw [line width=1.pt,dash pattern=on 2pt off 2pt,color=qqqqff] (0.32,3.64)-- (0.32,5.64);
\draw [line width=1.pt,dash pattern=on 2pt off 2pt,color=qqqqff] (0.32,3.64)-- (2.32,3.64);
\begin{scriptsize}
\draw [color=qqqqff] (-1.,3.)-- ++(-2.5pt,0 pt) -- ++(5.0pt,0 pt) ++(-2.5pt,-2.5pt) -- ++(0 pt,5.0pt);
\draw[color=qqqqff] (-0.9,2.83) node {$A$};
\draw [color=qqqqff] (1.,3.)-- ++(-2.5pt,0 pt) -- ++(5.0pt,0 pt) ++(-2.5pt,-2.5pt) -- ++(0 pt,5.0pt);
\draw[color=qqqqff] (0.98,2.75) node {$B$};
\draw [color=qqqqff] (1.,5.)-- ++(-2.5pt,0 pt) -- ++(5.0pt,0 pt) ++(-2.5pt,-2.5pt) -- ++(0 pt,5.0pt);
\draw[color=qqqqff] (1.14,5.37) node {$E$};
\draw [color=qqqqff] (-1.,5.)-- ++(-2.5pt,0 pt) -- ++(5.0pt,0 pt) ++(-2.5pt,-2.5pt) -- ++(0 pt,5.0pt);
\draw[color=qqqqff] (-0.86,5.37) node {$H$};
\draw [color=ududff] (2.32,3.64)-- ++(-2.5pt,0 pt) -- ++(5.0pt,0 pt) ++(-2.5pt,-2.5pt) -- ++(0 pt,5.0pt);
\draw[color=ududff] (2.46,4.01) node {$C$};
\draw [color=qqqqff] (2.32,5.64)-- ++(-2.5pt,0 pt) -- ++(5.0pt,0 pt) ++(-2.5pt,-2.5pt) -- ++(0 pt,5.0pt);
\draw[color=qqqqff] (2.46,6.01) node {$F$};
\draw [color=qqqqff] (0.32,5.64)-- ++(-2.5pt,0 pt) -- ++(5.0pt,0 pt) ++(-2.5pt,-2.5pt) -- ++(0 pt,5.0pt);
\draw[color=qqqqff] (0.46,6.01) node {$G$};
\draw [color=qqqqff] (0.32,3.64)-- ++(-2.5pt,0 pt) -- ++(5.0pt,0 pt) ++(-2.5pt,-2.5pt) -- ++(0 pt,5.0pt);
\draw[color=qqqqff] (0.46,4.01) node {$D$};
\end{scriptsize}
\end{tikzpicture}

\end{minipage}


\ExoCu{Représenter. Chercher.}

Un pavé droit a pour longueur $15$ cm, pour largeur $10$ cm et pour hauteur $20$ cm.

Calcule le volume du pavé droit en cm$^3$ et en dm$^3$.

\point{3}


\ExoCu{Calculer.}

\begin{minipage}{0.7\linewidth}

ABEH est un carré de coté de longueur 6 cm.  

Calcule le volume du cube ABCDEFGH en cm$^3$.

\point{4}

\end{minipage}
\hfill
\begin{minipage}{0.3\linewidth}

\definecolor{ududff}{rgb}{0.30196078431372547,0.30196078431372547,1.}
\definecolor{qqqqff}{rgb}{0.,0.,1.}
\definecolor{eqeqeq}{rgb}{0.8784313725490196,0.8784313725490196,0.8784313725490196}
\definecolor{qqttcc}{rgb}{0.,0.2,0.8}
\begin{tikzpicture}[line cap=round,line join=round,>=triangle 45,x=0.7462686567164183cm,y=0.7462686567164183cm]
\clip(-2.,1.52) rectangle (3.36,6.94);
\fill[line width=1.pt,color=eqeqeq,fill=eqeqeq,fill opacity=0.10000000149011612] (-1.76,1.86) -- (3.12,1.86) -- (3.12,6.74) -- (-1.76,6.74) -- cycle;
\draw [color=qqttcc](-0.72,3.02) node[anchor=north west] {$6\text{ cm}$};
\draw [line width=1.pt,color=eqeqeq] (-1.76,1.86)-- (3.12,1.86);
\draw [line width=1.pt,color=eqeqeq] (3.12,1.86)-- (3.12,6.74);
\draw [line width=1.pt,color=eqeqeq] (3.12,6.74)-- (-1.76,6.74);
\draw [line width=1.pt,color=eqeqeq] (-1.76,6.74)-- (-1.76,1.86);
\draw [line width=1.pt,color=eqeqeq] (-1.14,2.9)-- (0.58,2.9);
\draw [line width=1.pt,color=qqqqff] (-1.14,2.9)-- (0.58,2.9);
\draw [line width=1.pt,color=qqqqff] (0.58,2.9)-- (0.58,4.62);
\draw [line width=1.pt,color=qqqqff] (0.58,4.62)-- (-1.14,4.62);
\draw [line width=1.pt,color=qqqqff] (-1.14,4.62)-- (-1.14,2.9);
\draw [line width=1.pt,color=qqqqff] (-1.14,4.62)-- (0.94,5.56);
\draw [line width=1.pt,color=qqqqff] (0.94,5.56)-- (2.66,5.56);
\draw [line width=1.pt,color=qqqqff] (2.66,5.56)-- (2.66,3.84);
\draw [line width=1.pt,color=qqqqff] (2.66,3.84)-- (0.58,2.9);
\draw [line width=1.pt,color=qqqqff] (0.58,4.62)-- (2.66,5.56);
\draw [line width=1.pt,dash pattern=on 2pt off 2pt,color=qqqqff] (-1.14,2.9)-- (0.94,3.84);
\draw [line width=1.pt,dash pattern=on 2pt off 2pt,color=qqqqff] (0.94,3.84)-- (0.94,5.56);
\draw [line width=1.pt,dash pattern=on 2pt off 2pt,color=qqqqff] (0.94,3.84)-- (2.66,3.84);
\draw [color=qqttcc](1.42,3.48) node[anchor=north west] {$11 \text{ cm}$};
\begin{scriptsize}
\draw [color=qqqqff] (-1.14,2.9)-- ++(-2.5pt,0 pt) -- ++(5.0pt,0 pt) ++(-2.5pt,-2.5pt) -- ++(0 pt,5.0pt);
\draw[color=qqqqff] (-1.04,2.73) node {$A$};
\draw [color=qqqqff] (0.58,2.9)-- ++(-2.5pt,0 pt) -- ++(5.0pt,0 pt) ++(-2.5pt,-2.5pt) -- ++(0 pt,5.0pt);
\draw[color=qqqqff] (0.56,2.65) node {$B$};
\draw [color=qqqqff] (0.58,4.62)-- ++(-2.5pt,0 pt) -- ++(5.0pt,0 pt) ++(-2.5pt,-2.5pt) -- ++(0 pt,5.0pt);
\draw[color=qqqqff] (0.72,4.99) node {$E$};
\draw [color=qqqqff] (-1.14,4.62)-- ++(-2.5pt,0 pt) -- ++(5.0pt,0 pt) ++(-2.5pt,-2.5pt) -- ++(0 pt,5.0pt);
\draw[color=qqqqff] (-1.,4.99) node {$H$};
\draw [color=ududff] (2.66,3.84)-- ++(-2.5pt,0 pt) -- ++(5.0pt,0 pt) ++(-2.5pt,-2.5pt) -- ++(0 pt,5.0pt);
\draw[color=ududff] (2.8,4.21) node {$C$};
\draw [color=qqqqff] (2.66,5.56)-- ++(-2.5pt,0 pt) -- ++(5.0pt,0 pt) ++(-2.5pt,-2.5pt) -- ++(0 pt,5.0pt);
\draw[color=qqqqff] (2.8,5.93) node {$F$};
\draw [color=qqqqff] (0.94,5.56)-- ++(-2.5pt,0 pt) -- ++(5.0pt,0 pt) ++(-2.5pt,-2.5pt) -- ++(0 pt,5.0pt);
\draw[color=qqqqff] (1.08,5.93) node {$G$};
\draw [color=qqqqff] (0.94,3.84)-- ++(-2.5pt,0 pt) -- ++(5.0pt,0 pt) ++(-2.5pt,-2.5pt) -- ++(0 pt,5.0pt);
\draw[color=qqqqff] (1.08,4.21) node {$D$};
\end{scriptsize}
\end{tikzpicture}

\end{minipage}






\subsection{ Utiliser les unités de volume }
\subsection{ Relier les unités de volume et de contenance (1 L = 1 dm3)}
\subsection{ Estimer si un angle est droit, aigu ou obtus}
\subsection{ Utiliser un rapporteur pour mesurer un angle en degrés}
\subsection{ Construire à l'aide du rapporteur, un angle de mesure donnée en degrés}
\subsection{ Comparer, estimer, mesurer des longueurs et des périmètres}



\end{document}
